
\documentclass[10pt]{article}
\usepackage[T1]{fontenc}
\usepackage{lmodern}
\usepackage[margin=0.72in]{geometry}
\usepackage{microtype}
\usepackage{amsmath,amssymb,amsthm}
\usepackage{booktabs,longtable,array,ragged2e}
\usepackage{enumitem}
\usepackage{xcolor}
\usepackage{hyperref}
\usepackage{xurl}
\usepackage{fancyhdr}
\usepackage{titlesec}
\usepackage{pdflscape}
\usepackage{parskip}

\definecolor{navy}{HTML}{17365D}
\definecolor{blue}{HTML}{1F4E79}
\definecolor{teal}{HTML}{0F6B78}
\definecolor{graytext}{HTML}{4A535A}
\definecolor{lightblue}{HTML}{EEF5FA}
\definecolor{lightgray}{HTML}{F5F6F7}
\definecolor{darkred}{HTML}{7A2430}

\hypersetup{
  colorlinks=true,
  linkcolor=navy,
  urlcolor=blue,
  citecolor=teal
}
\urlstyle{same}
\setlength{\emergencystretch}{3em}
\setlength{\parindent}{0pt}

\pagestyle{fancy}
\fancyhf{}
\lhead{\textcolor{navy}{Mathematical Technique Proof Library}}
\rhead{\textcolor{graytext}{Proof dependency architecture}}
\cfoot{\thepage}
\setlength{\headheight}{14pt}

\titleformat{\section}{\Large\bfseries\color{navy}}{\thesection}{0.7em}{}
\titleformat{\subsection}{\large\bfseries\color{blue}}{\thesubsection}{0.6em}{}
\titleformat{\subsubsection}{\normalsize\bfseries\color{teal}}{\thesubsubsection}{0.5em}{}

\newcolumntype{L}[1]{>{\raggedright\arraybackslash\hspace{0pt}}p{#1}}
\newcolumntype{C}[1]{>{\centering\arraybackslash\hspace{0pt}}p{#1}}

\newcommand{\pidraw}[1]{\begingroup\urlstyle{tt}\nolinkurl{#1}\endgroup}
\newcommand{\pidref}[1]{\hyperref[prf:#1]{\pidraw{#1}}}
\newcommand{\pid}[1]{\pidref{#1}}
\newcommand{\nr}{\emph{Numerical Recipes: The Art of Scientific Computing, Third Edition}}

\newtheorem{theorem}{Theorem}[section]
\newtheorem{lemma}[theorem]{Lemma}
\newtheorem{proposition}[theorem]{Proposition}
\newtheorem{corollary}[theorem]{Corollary}
\theoremstyle{definition}
\newtheorem{definition}[theorem]{Definition}
\theoremstyle{remark}
\newtheorem*{remark}{Remark}

\newenvironment{proofstrategy}
{\begin{quote}\small\noindent\textbf{\textcolor{teal}{Proof Strategy.}}\ }
{\end{quote}}

\newenvironment{proofanalysis}
{\begin{quote}\small\noindent\textbf{\textcolor{blue}{Proof Analysis.}}\ }
{\end{quote}}

\newenvironment{computationalrelevance}
{\begin{quote}\small\noindent\textbf{\textcolor{navy}{Computational Relevance.}}\ }
{\end{quote}}

\newenvironment{scopelimitation}
{\begin{quote}\small\noindent\textbf{\textcolor{darkred}{Scope Limitation.}}\ }
{\end{quote}}

\newcommand{\meta}[2]{\noindent\textbf{#1:} #2\par}
\hypersetup{pdftitle={Mathematical Technique Proof Library - Modular Referenced Proofs}}
\begin{document}
\hypersetup{pageanchor=false}

\begin{titlepage}
\centering
\vspace*{1.1in}
{\Huge\bfseries\color{navy} Mathematical Technique Proof Library\par}
\vspace{0.25in}
{\Large\color{blue} Version A --- Modular Referenced Proofs\par}
\vspace{0.5in}
{\large A dependency-aware proof architecture for calculus and scientific computing\par}
\vfill
\begin{minipage}{0.84\textwidth}
\small
\textbf{Purpose.} This library classifies the mathematical techniques in the supplied
computational catalog and proves the statements that can be rigorously supported by the
supplied sources. Proof presentation follows the methodology of
\emph{Mathematical Proofs: A Transition to Advanced Mathematics}, Fourth Edition.
Numerical context is traced to \nr. The calculus source is referenced generically as the
``source calculus framework.''

\medskip
\textbf{Boundary.} Definitions are stated rather than proved. Deep results whose complete
proof would require theory outside the supplied proof sources are explicitly marked as
source-bound foundational results or scope limitations.
\end{minipage}
\vfill
{\large September 23, 2026\par}
\end{titlepage}

\clearpage
\hypersetup{pageanchor=true}
\pagenumbering{roman}
\tableofcontents
\clearpage
\pagenumbering{arabic}

\section{Architecture and Conventions}

The proof system distinguishes mathematical objects from algorithms and software implementations.
A calculus concept such as a derivative is a mathematical object; a centered finite difference is a
numerical approximation; and a particular library routine is software. Only mathematical statements
receive proofs.

The dependency relation is directed from prerequisite to dependent result. Every proof object has a
stable proof ID. The modular document may invoke earlier IDs. The complete document expands every
non-foundational dependency locally before the target proof.

\subsection{Proof-method vocabulary}

The proof-method terminology follows \emph{Mathematical Proofs}: direct proof, proof by
contrapositive, proof by cases, contradiction, existence proof, induction, and proof evaluation.
Epsilon-delta arguments are treated as direct proofs of quantified implications. Numerical
convergence arguments may combine direct proof, induction, inequalities, invariants, and limits.

\subsection{Foundational boundary}

The library does not attempt to derive all mathematics from set theory. Expansion stops at explicitly
declared foundations: classical logic, the ordered-field axioms, completeness of the real numbers,
selected compact-interval and Riemann-integral facts, elementary probability identities, basic
finite-dimensional linear algebra, and complex-exponential identities. This boundary is visible
rather than implicit.
\section{Definitions and Foundational Assumptions}

\subsection{Foundations}

\phantomsection\label{prf:AX-CLASSICAL-LOGIC}

\paragraph{\pidraw{AX-CLASSICAL-LOGIC} --- Classical logic and quantified implication}

Classical propositional logic, quantified logic, equality substitution, and ordinary rules of inference are accepted as the logical foundation.

\smallskip\noindent\textit{Traceability:} \textbf{Proof method:} Mathematical Proofs, Ch. 2\par

\phantomsection\label{prf:AX-COMPLEX-EXP}

\paragraph{\pidraw{AX-COMPLEX-EXP} --- Complex exponential identities}

Euler's formula and the usual laws of complex exponentials are accepted, including $e^{u+v}=e^u e^v$.

\smallskip\noindent\textit{Traceability:} \textbf{Proof method:} Mathematical Proofs, online Ch. 18.5--18.6\par

\phantomsection\label{prf:AX-ELEMENTARY-SEQUENCE-LIMITS}

\paragraph{\pidraw{AX-ELEMENTARY-SEQUENCE-LIMITS} --- Elementary sequence-limit facts}

The elementary sequence limit $r^n\to0$ for $|r|<1$, stability under index shifts, and basic real-sequence limit algebra are accepted.

\smallskip\noindent\textit{Traceability:} \textbf{Proof method:} Mathematical Proofs, Sec. 14.1\par

\phantomsection\label{prf:AX-EVT}

\paragraph{\pidraw{AX-EVT} --- Extreme Value Theorem}

If $f:[a,b]\to\mathbb{R}$ is continuous, then $f$ attains an absolute maximum and an absolute minimum on $[a,b]$.

\begin{scopelimitation}Accepted at the foundational boundary because a full compactness proof is beyond the proof development supplied in the core printed chapters.\end{scopelimitation}

\smallskip\noindent\textit{Traceability:} \textbf{Mathematical source:} source calculus framework, Sec. 3.1\par

\phantomsection\label{prf:AX-LINEAR-ALGEBRA-BASICS}

\paragraph{\pidraw{AX-LINEAR-ALGEBRA-BASICS} --- Basic finite-dimensional linear algebra}

The standard algebra of real vectors and matrices is accepted: transpose identities, inner-product bilinearity, and the equivalence between nonsingularity and uniqueness of solutions of homogeneous linear systems.

\smallskip\noindent\textit{Traceability:} \textbf{Proof method:} Mathematical Proofs, online Ch. 17\par

\phantomsection\label{prf:AX-POLY-FACTOR}

\paragraph{\pidraw{AX-POLY-FACTOR} --- Polynomial factor theorem}

If a polynomial $p$ satisfies $p(r)=0$, then $p(x)=(x-r)q(x)$ for some polynomial $q$ of degree one less than $p$ when $p$ is nonzero.

\smallskip\noindent\textit{Traceability:} \emph{No additional source locator required.}\par

\phantomsection\label{prf:AX-PROBABILITY}

\paragraph{\pidraw{AX-PROBABILITY} --- Linearity of expectation and variance for independent variables}

Expectation is linear. For independent random variables with finite variances, variances add. These are the probabilistic foundations used for the Monte Carlo results.

\begin{scopelimitation}The supplied proof-method text does not develop probability theory; these probability identities terminate proof expansion.\end{scopelimitation}

\smallskip\noindent\textit{Traceability:} \textbf{Numerical source:} Numerical Recipes, Chs. 7 and 14\par

\phantomsection\label{prf:AX-REAL-COMPLETENESS}

\paragraph{\pidraw{AX-REAL-COMPLETENESS} --- Completeness of the real numbers}

Every nonempty subset of $\mathbb{R}$ that is bounded above has a least upper bound in $\mathbb{R}$.

\smallskip\noindent\textit{Traceability:} \textbf{Proof method:} Mathematical Proofs, online Ch. 18.2\par

\phantomsection\label{prf:AX-REAL-FIELD}

\paragraph{\pidraw{AX-REAL-FIELD} --- Ordered-field properties of the real numbers}

The real numbers form an ordered field. In particular, the algebraic and order properties of $+$, multiplication, absolute value, and inequalities are available.

\smallskip\noindent\textit{Traceability:} \textbf{Proof method:} Mathematical Proofs, Ch. 4.3; online Ch. 18.1\par

\phantomsection\label{prf:AX-RIEMANN-BASIC}

\paragraph{\pidraw{AX-RIEMANN-BASIC} --- Basic Riemann-integral order and integrability facts}

Continuous functions on closed bounded intervals are Riemann integrable; the Riemann integral is linear and order preserving; and if $m\le f\le M$ on $[a,b]$, then $m(b-a)\le\int_a^b f\le M(b-a)$.

\smallskip\noindent\textit{Traceability:} \textbf{Mathematical source:} source calculus framework, Secs. 4.3--4.4\par

\phantomsection\label{prf:DEF-BIG-O}

\paragraph{\pidraw{DEF-BIG-O} --- Big-O remainder notation}

The statement $R(h)=O(h^p)$ as $h\to0$ means that there exist constants $C,\eta>0$ such that $|R(h)|\le C|h|^p$ whenever $0<|h|<\eta$.

\smallskip\noindent\textit{Traceability:} \textbf{Numerical source:} Numerical Recipes, Sec. 1.1\par

\phantomsection\label{prf:DEF-COMPOSITION}

\paragraph{\pidraw{DEF-COMPOSITION} --- Composition of functions}

For functions $g:X\to Y$ and $f:Y\to Z$, the composition is $(f\circ g)(x)=f(g(x))$.

\smallskip\noindent\textit{Traceability:} \textbf{Proof method:} Mathematical Proofs, Sec. 10.4; \textbf{Mathematical source:} source calculus framework, function review\par

\phantomsection\label{prf:DEF-CONTINUITY}

\paragraph{\pidraw{DEF-CONTINUITY} --- Continuity at a point}

A function $f$ is continuous at $a$ when $f(a)$ is defined and $\lim_{x\to a}f(x)=f(a)$.

\smallskip\noindent\textit{Traceability:} \textbf{Proof method:} Mathematical Proofs, Sec. 14.5; \textbf{Mathematical source:} source calculus framework, Sec. 1.4\par

\phantomsection\label{prf:DEF-DERIVATIVE}

\paragraph{\pidraw{DEF-DERIVATIVE} --- Derivative at a point}

The derivative of $f$ at $a$ is
\[
f'(a)=\lim_{x\to a}\frac{f(x)-f(a)}{x-a},
\]
provided the limit exists.

\smallskip\noindent\textit{Traceability:} \textbf{Proof method:} Mathematical Proofs, Sec. 14.6; \textbf{Mathematical source:} source calculus framework, Sec. 2.1\par

\phantomsection\label{prf:DEF-FUNCTION}

\paragraph{\pidraw{DEF-FUNCTION} --- Function}

A function $f:X\to Y$ assigns to each $x\in X$ exactly one value $f(x)\in Y$.

\smallskip\noindent\textit{Traceability:} \textbf{Proof method:} Mathematical Proofs, Sec. 10.1; \textbf{Mathematical source:} source calculus framework, Sec. P.3\par

\phantomsection\label{prf:DEF-INVERSE-FUNCTION}

\paragraph{\pidraw{DEF-INVERSE-FUNCTION} --- Inverse function}

A bijection $f:X\to Y$ has inverse $f^{-1}:Y\to X$ defined by $f^{-1}(f(x))=x$ and $f(f^{-1}(y))=y$.

\smallskip\noindent\textit{Traceability:} \textbf{Proof method:} Mathematical Proofs, Sec. 10.5; \textbf{Mathematical source:} source calculus framework, Sec. 5.3\par

\phantomsection\label{prf:DEF-LIMIT-FUNCTION}

\paragraph{\pidraw{DEF-LIMIT-FUNCTION} --- Limit of a real function}

For a function $f$ defined on a deleted neighborhood of $a$, $\lim_{x\to a}f(x)=L$ means that for every $\varepsilon>0$ there exists $\delta>0$ such that
\[
0<|x-a|<\delta\quad\Longrightarrow\quad |f(x)-L|<\varepsilon.
\]

\smallskip\noindent\textit{Traceability:} \textbf{Proof method:} Mathematical Proofs, Sec. 14.3; \textbf{Mathematical source:} source calculus framework, Secs. 1.2--1.4\par

\phantomsection\label{prf:DEF-RIEMANN-INTEGRAL}

\paragraph{\pidraw{DEF-RIEMANN-INTEGRAL} --- Definite Riemann integral}

The definite integral $\int_a^b f(x)\,dx$ is the common limit of Riemann sums when that limit exists independently of the tagged partitions whose mesh tends to zero.

\smallskip\noindent\textit{Traceability:} \textbf{Mathematical source:} source calculus framework, Sec. 4.3\par

\phantomsection\label{prf:DEF-SERIES}

\paragraph{\pidraw{DEF-SERIES} --- Convergence of an infinite series}

For a series $\sum_{k=1}^{\infty}a_k$, let $s_n=\sum_{k=1}^n a_k$. The series converges to $s$ when $s_n\to s$.

\smallskip\noindent\textit{Traceability:} \textbf{Proof method:} Mathematical Proofs, Sec. 14.2; \textbf{Mathematical source:} source calculus framework, Sec. 9.2\par

\subsection{Optimization}

\phantomsection\label{prf:DEF-CRITICAL-POINT}

\paragraph{\pidraw{DEF-CRITICAL-POINT} --- Critical point}

An interior point $c$ of the domain is a critical point of $f$ when $f'(c)=0$ or $f'(c)$ does not exist.

\smallskip\noindent\textit{Traceability:} \textbf{Mathematical source:} source calculus framework, Sec. 3.1\par

\subsection{Integral Calculus}

\phantomsection\label{prf:DEF-RIEMANN-SUM}

\paragraph{\pidraw{DEF-RIEMANN-SUM} --- Riemann sum}

For a tagged partition $a=x_0<\cdots<x_n=b$ with tags $c_i\in[x_{i-1},x_i]$, a Riemann sum is
 \[
 \sum_{i=1}^n f(c_i)\Delta x_i,\qquad \Delta x_i=x_i-x_{i-1}.
 \]

\smallskip\noindent\textit{Traceability:} \textbf{Mathematical source:} source calculus framework, Sec. 4.3\par

\subsection{Multivariable Calculus}

\phantomsection\label{prf:DEF-DIFFERENTIABLE-RN}

\paragraph{\pidraw{DEF-DIFFERENTIABLE-RN} --- Differentiability in several variables}

A function $F:\mathbb{R}^n\to\mathbb{R}$ is differentiable at $a$ if there is a linear map $L$ such that
\[
F(a+h)=F(a)+L(h)+r(h),\qquad \frac{r(h)}{\|h\|}\to0.
\]
When partial derivatives identify $L$, $L(h)=\nabla F(a)\cdot h$.

\smallskip\noindent\textit{Traceability:} \textbf{Mathematical source:} source calculus framework, Secs. 13.3--13.6\par

\phantomsection\label{prf:DEF-DIRECTIONAL-DERIVATIVE}

\paragraph{\pidraw{DEF-DIRECTIONAL-DERIVATIVE} --- Directional derivative}

For a unit vector $u$, the directional derivative of $F$ at $a$ in direction $u$ is
\[
D_uF(a)=\lim_{t\to0}\frac{F(a+tu)-F(a)}{t},
\]
when the limit exists.

\smallskip\noindent\textit{Traceability:} \textbf{Mathematical source:} source calculus framework, Sec. 13.6\par

\phantomsection\label{prf:DEF-GRADIENT}

\paragraph{\pidraw{DEF-GRADIENT} --- Gradient}

When the first partial derivatives exist, the gradient is
 \[
 \nabla F=(F_{x_1},\ldots,F_{x_n}).
 \]

\smallskip\noindent\textit{Traceability:} \textbf{Mathematical source:} source calculus framework, Sec. 13.6\par

\phantomsection\label{prf:DEF-PARTIAL-DERIVATIVE}

\paragraph{\pidraw{DEF-PARTIAL-DERIVATIVE} --- Partial derivative}

For $F:\mathbb{R}^n\to\mathbb{R}$, the $i$th partial derivative at $a$ is the one-variable derivative obtained by varying only the $i$th coordinate.

\smallskip\noindent\textit{Traceability:} \textbf{Mathematical source:} source calculus framework, Sec. 13.3\par

\subsection{Vector Analysis}

\phantomsection\label{prf:DEF-LINE-INTEGRAL}

\paragraph{\pidraw{DEF-LINE-INTEGRAL} --- Line integral of a vector field}

For a piecewise smooth parametrized curve $r:[\alpha,\beta]\to\mathbb{R}^n$,
 \[
 \int_C F\cdot dr=\int_\alpha^\beta F(r(t))\cdot r'(t)\,dt.
 \]

\smallskip\noindent\textit{Traceability:} \textbf{Mathematical source:} source calculus framework, Sec. 15.2\par

\phantomsection\label{prf:DEF-VECTOR-FIELD}

\paragraph{\pidraw{DEF-VECTOR-FIELD} --- Vector field}

A vector field assigns a vector $F(x)$ to each point $x$ in its domain.

\smallskip\noindent\textit{Traceability:} \textbf{Mathematical source:} source calculus framework, Sec. 15.1\par

\subsection{Numerical Approximation}

\phantomsection\label{prf:DEF-INTERPOLATION}

\paragraph{\pidraw{DEF-INTERPOLATION} --- Interpolation problem}

Interpolation seeks a function from a chosen class that exactly matches prescribed data values at specified nodes.

\smallskip\noindent\textit{Traceability:} \emph{No additional source locator required.}\par

\subsection{Linear Algebra}

\phantomsection\label{prf:DEF-EIGENPAIR}

\paragraph{\pidraw{DEF-EIGENPAIR} --- Eigenvalue and eigenvector}

For a square matrix $A$, a nonzero vector $v$ is an eigenvector with eigenvalue $\lambda$ when $Av=\lambda v$.

\smallskip\noindent\textit{Traceability:} \emph{No additional source locator required.}\par

\subsection{Differential Equations}

\phantomsection\label{prf:DEF-ODE-IVP}

\paragraph{\pidraw{DEF-ODE-IVP} --- Initial-value problem for an ODE}

An initial-value problem prescribes a differential equation $y'=f(t,y)$ together with an initial condition $y(t_0)=y_0$.

\smallskip\noindent\textit{Traceability:} \textbf{Mathematical source:} source calculus framework, Sec. 6.1\par

\subsection{Spectral Methods}

\phantomsection\label{prf:DEF-DFT}

\paragraph{\pidraw{DEF-DFT} --- Discrete Fourier transform}

For $N$ samples $x_0,\ldots,x_{N-1}$, the DFT is
 \[
 X_k=\sum_{n=0}^{N-1}x_n e^{-2\pi i kn/N},\qquad 0\le k<N.
 \]

\smallskip\noindent\textit{Traceability:} \emph{No additional source locator required.}\par\section{Source-Bound Results and Non-Proof Objects}

The following catalog techniques are intentionally not assigned synthetic proofs. A source-bound result is mathematically meaningful but its complete proof requires theory beyond the selected source boundary. A non-proof object is a workflow, implementation family, or collection of definitions rather than one proposition.

\subsection{Foundations}

\phantomsection\label{prf:NOP-MODELING}

\paragraph{\pidraw{NOP-MODELING} --- Mathematical modeling as a workflow}

Model construction is a representational workflow: assumptions, variables, equations, and validation. Individual claims inside a model may be proved, but modeling itself is not one proposition.

\smallskip\noindent\textit{Traceability:} \textbf{Mathematical source:} source calculus framework, Secs. P.1--P.3\par

\subsection{Multivariable Calculus}

\phantomsection\label{prf:ADV-LAGRANGE}

\paragraph{\pidraw{ADV-LAGRANGE} --- Lagrange multiplier theorem}

A rigorous general proof normally invokes an implicit-function/rank theorem. The source calculus framework states and applies the method, but the supplied proof-method source does not develop that theorem.

\begin{scopelimitation}A rigorous general proof normally invokes an implicit-function/rank theorem. The source calculus framework states and applies the method, but the supplied proof-method source does not develop that theorem.\end{scopelimitation}

\smallskip\noindent\textit{Traceability:} \textbf{Mathematical source:} source calculus framework, Sec. 13.10\par

\subsection{Vector Analysis}

\phantomsection\label{prf:ADV-DIVERGENCE}

\paragraph{\pidraw{ADV-DIVERGENCE} --- Divergence theorem}

The theorem is retained as a source-bound foundational result; a full proof requires geometric measure/partition arguments beyond the selected proof boundary.

\begin{scopelimitation}The theorem is retained as a source-bound foundational result; a full proof requires geometric measure/partition arguments beyond the selected proof boundary.\end{scopelimitation}

\smallskip\noindent\textit{Traceability:} \textbf{Mathematical source:} source calculus framework, Sec. 15.7\par

\phantomsection\label{prf:ADV-GREEN}

\paragraph{\pidraw{ADV-GREEN} --- Green's theorem}

The theorem is retained as a source-bound foundational result for this library; a full proof requires planar integration theory not developed in the supplied proof-method source.

\begin{scopelimitation}The theorem is retained as a source-bound foundational result for this library; a full proof requires planar integration theory not developed in the supplied proof-method source.\end{scopelimitation}

\smallskip\noindent\textit{Traceability:} \textbf{Mathematical source:} source calculus framework, Sec. 15.4\par

\phantomsection\label{prf:ADV-STOKES}

\paragraph{\pidraw{ADV-STOKES} --- Stokes's theorem}

The theorem is retained as a source-bound foundational result; a full proof requires surface orientation, parametrization, and patching arguments beyond the selected proof boundary.

\begin{scopelimitation}The theorem is retained as a source-bound foundational result; a full proof requires surface orientation, parametrization, and patching arguments beyond the selected proof boundary.\end{scopelimitation}

\smallskip\noindent\textit{Traceability:} \textbf{Mathematical source:} source calculus framework, Sec. 15.8\par

\subsection{Numerical Approximation}

\phantomsection\label{prf:ADV-CUBIC-SPLINE}

\paragraph{\pidraw{ADV-CUBIC-SPLINE} --- Existence/uniqueness and error theory for cubic splines}

A full natural/not-a-knot spline existence and error theory requires a longer linear-algebra and approximation-theory development than the supplied core proof source provides.

\begin{scopelimitation}A full natural/not-a-knot spline existence and error theory requires a longer linear-algebra and approximation-theory development than the supplied core proof source provides.\end{scopelimitation}

\smallskip\noindent\textit{Traceability:} \textbf{Numerical source:} Numerical Recipes, Sec. 3.3\par

\phantomsection\label{prf:NOP-SPECIAL-FUNCTIONS}

\paragraph{\pidraw{NOP-SPECIAL-FUNCTIONS} --- Special-function library evaluations}

Special functions are mathematical objects with many identities, but the catalog entry denotes an implementation family rather than one proposition. No single proof is manufactured for it.

\smallskip\noindent\textit{Traceability:} \textbf{Numerical source:} Numerical Recipes, Ch. 6\par

\subsection{Root Finding}

\phantomsection\label{prf:ADV-BRENT-ROOT}

\paragraph{\pidraw{ADV-BRENT-ROOT} --- Full convergence theory for safeguarded Brent-Dekker root finding}

The sign-bracketing invariant is proved in this library, but a full rate/convergence proof for the safeguarded interpolation logic is outside the source-supported proof core.

\begin{scopelimitation}The sign-bracketing invariant is proved in this library, but a full rate/convergence proof for the safeguarded interpolation logic is outside the source-supported proof core.\end{scopelimitation}

\smallskip\noindent\textit{Traceability:} \textbf{Numerical source:} Numerical Recipes, Sec. 9.3\par

\phantomsection\label{prf:ADV-MULTIROOT}

\paragraph{\pidraw{ADV-MULTIROOT} --- Local convergence theory for Newton methods on nonlinear systems}

A rigorous theorem requires Jacobian invertibility plus multivariable remainder estimates; the catalog records the method family but the supplied proof source does not develop this numerical-analysis theorem.

\begin{scopelimitation}A rigorous theorem requires Jacobian invertibility plus multivariable remainder estimates; the catalog records the method family but the supplied proof source does not develop this numerical-analysis theorem.\end{scopelimitation}

\smallskip\noindent\textit{Traceability:} \textbf{Numerical source:} Numerical Recipes, Secs. 9.6--9.7\par

\subsection{Optimization Algorithms}

\phantomsection\label{prf:ADV-BFGS}

\paragraph{\pidraw{ADV-BFGS} --- Convergence theory for BFGS/quasi-Newton methods}

Rigorous global/local convergence theorems require assumptions and line-search theory not supplied in the proof-method source.

\begin{scopelimitation}Rigorous global/local convergence theorems require assumptions and line-search theory not supplied in the proof-method source.\end{scopelimitation}

\smallskip\noindent\textit{Traceability:} \textbf{Numerical source:} Numerical Recipes, Sec. 10.9\par

\phantomsection\label{prf:ADV-BRENT-MIN}

\paragraph{\pidraw{ADV-BRENT-MIN} --- Convergence guarantees for Brent minimization}

The hybrid algorithm combines safeguarded interpolation and golden-section steps; a complete algorithmic convergence proof is outside the source-supported proof core.

\begin{scopelimitation}The hybrid algorithm combines safeguarded interpolation and golden-section steps; a complete algorithmic convergence proof is outside the source-supported proof core.\end{scopelimitation}

\smallskip\noindent\textit{Traceability:} \textbf{Numerical source:} Numerical Recipes, Sec. 10.3\par

\phantomsection\label{prf:ADV-NELDER-MEAD}

\paragraph{\pidraw{ADV-NELDER-MEAD} --- Convergence of the multidimensional simplex method}

General convergence behavior is subtle and is not established by the supplied sources as a simple calculus theorem.

\begin{scopelimitation}General convergence behavior is subtle and is not established by the supplied sources as a simple calculus theorem.\end{scopelimitation}

\smallskip\noindent\textit{Traceability:} \textbf{Numerical source:} Numerical Recipes, Sec. 10.5\par

\phantomsection\label{prf:ADV-SIMAN}

\paragraph{\pidraw{ADV-SIMAN} --- Convergence of simulated annealing}

General asymptotic convergence depends on stochastic-process and cooling-schedule theory outside the supplied proof-method source.

\begin{scopelimitation}General asymptotic convergence depends on stochastic-process and cooling-schedule theory outside the supplied proof-method source.\end{scopelimitation}

\smallskip\noindent\textit{Traceability:} \textbf{Numerical source:} Numerical Recipes, Sec. 10.12\par

\subsection{Numerical Integration}

\phantomsection\label{prf:ADV-GAUSSIAN-QUADRATURE}

\paragraph{\pidraw{ADV-GAUSSIAN-QUADRATURE} --- Degree-\$2n-1\$ exactness of Gaussian quadrature}

The standard proof requires orthogonal-polynomial theory and a polynomial-division argument not developed in the supplied proof-method source.

\begin{scopelimitation}The standard proof requires orthogonal-polynomial theory and a polynomial-division argument not developed in the supplied proof-method source.\end{scopelimitation}

\smallskip\noindent\textit{Traceability:} \textbf{Numerical source:} Numerical Recipes, Sec. 4.6\par

\phantomsection\label{prf:ADV-LEBEDEV}

\paragraph{\pidraw{ADV-LEBEDEV} --- Exactness theory for Lebedev spherical quadrature}

The catalog uses spherical quadrature as an implementation technique, but exactness/design theory for Lebedev rules is outside the supplied proof sources.

\begin{scopelimitation}The catalog uses spherical quadrature as an implementation technique, but exactness/design theory for Lebedev rules is outside the supplied proof sources.\end{scopelimitation}

\smallskip\noindent\textit{Traceability:} \emph{No additional source locator required.}\par

\subsection{Linear Algebra}

\phantomsection\label{prf:ADV-CG}

\paragraph{\pidraw{ADV-CG} --- Finite-termination/convergence theory for conjugate gradients}

A full proof requires positive-definite quadratic forms, Krylov spaces, and conjugacy theory beyond the supplied proof-method development.

\begin{scopelimitation}A full proof requires positive-definite quadratic forms, Krylov spaces, and conjugacy theory beyond the supplied proof-method development.\end{scopelimitation}

\smallskip\noindent\textit{Traceability:} \textbf{Numerical source:} Numerical Recipes, Sec. 10.8\par

\phantomsection\label{prf:ADV-SVD}

\paragraph{\pidraw{ADV-SVD} --- Existence and numerical theory of the singular value decomposition}

The algebraic existence proof and numerical perturbation theory require a broader linear-algebra foundation than the printed proof-method chapters.

\begin{scopelimitation}The algebraic existence proof and numerical perturbation theory require a broader linear-algebra foundation than the printed proof-method chapters.\end{scopelimitation}

\smallskip\noindent\textit{Traceability:} \textbf{Numerical source:} Numerical Recipes, Sec. 2.6\par

\subsection{Differential Equations}

\phantomsection\label{prf:ADV-BVP}

\paragraph{\pidraw{ADV-BVP} --- General convergence of shooting/relaxation for boundary-value problems}

General BVP convergence depends on problem-specific existence, conditioning, and nonlinear-solver hypotheses.

\begin{scopelimitation}General BVP convergence depends on problem-specific existence, conditioning, and nonlinear-solver hypotheses.\end{scopelimitation}

\smallskip\noindent\textit{Traceability:} \textbf{Numerical source:} Numerical Recipes, Ch. 18\par

\phantomsection\label{prf:ADV-RUNGE-KUTTA}

\paragraph{\pidraw{ADV-RUNGE-KUTTA} --- Order and convergence theory for Runge-Kutta methods}

The catalog's production methods require Butcher-order conditions and stability theory beyond the supplied proof-method source. The library proves Euler local truncation and Richardson cancellation as supported foundations.

\begin{scopelimitation}The catalog's production methods require Butcher-order conditions and stability theory beyond the supplied proof-method source. The library proves Euler local truncation and Richardson cancellation as supported foundations.\end{scopelimitation}

\smallskip\noindent\textit{Traceability:} \textbf{Numerical source:} Numerical Recipes, Secs. 17.1--17.4\par

\phantomsection\label{prf:ADV-STIFF}

\paragraph{\pidraw{ADV-STIFF} --- Stability theory for stiff ODE solvers}

A-stability, L-stability, multistep stability regions, and stiff convergence require a dedicated numerical ODE theory.

\begin{scopelimitation}A-stability, L-stability, multistep stability regions, and stiff convergence require a dedicated numerical ODE theory.\end{scopelimitation}

\smallskip\noindent\textit{Traceability:} \textbf{Numerical source:} Numerical Recipes, Secs. 17.5--17.6\par

\subsection{Spectral Methods}

\phantomsection\label{prf:ADV-WAVELET}

\paragraph{\pidraw{ADV-WAVELET} --- General wavelet transform reconstruction theory}

A complete proof requires multiresolution and filter-bank theory not developed in the supplied sources.

\begin{scopelimitation}A complete proof requires multiresolution and filter-bank theory not developed in the supplied sources.\end{scopelimitation}

\smallskip\noindent\textit{Traceability:} \textbf{Numerical source:} Numerical Recipes, Sec. 13.10\par

\subsection{Probability and Monte Carlo}

\phantomsection\label{prf:NOP-DESCRIPTIVE-STATS}

\paragraph{\pidraw{NOP-DESCRIPTIVE-STATS} --- Descriptive statistics as a collection of definitions}

Means, variances, correlations, and summaries form a collection of definitions and estimators rather than one theorem. Specific estimator properties are proved separately when supported.

\smallskip\noindent\textit{Traceability:} \textbf{Numerical source:} Numerical Recipes, Ch. 14\par

\phantomsection\label{prf:NOP-RNG}

\paragraph{\pidraw{NOP-RNG} --- Pseudorandom generator implementation quality}

A software RNG is an implementation object. Mathematical properties of a distribution or estimator may be proved, but a particular library generator is not itself a theorem.

\smallskip\noindent\textit{Traceability:} \textbf{Numerical source:} Numerical Recipes, Secs. 7.1--7.6\par

\subsection{Geometry}

\phantomsection\label{prf:NOP-COMPUTATIONAL-GEOMETRY}

\paragraph{\pidraw{NOP-COMPUTATIONAL-GEOMETRY} --- Computational-geometry module family}

Computational geometry is a family of algorithms and predicates, not one theorem. Individual algorithms require their own specifications and correctness proofs.

\smallskip\noindent\textit{Traceability:} \textbf{Numerical source:} Numerical Recipes, Ch. 21\par
\section{Pilot Validation}
Before the full library was emitted, ten representative results spanning limits, derivatives,
optimization, integration, series, root finding, interpolation, numerical differentiation, ODE
methods, and multivariable calculus were checked for statement consistency, dependency closure, and
topological ordering.

\begin{center}
\small
\begin{tabular}{@{}L{0.31\textwidth}L{0.22\textwidth}C{0.11\textwidth}C{0.11\textwidth}C{0.11\textwidth}@{}}
\toprule
\textbf{Proof ID} & \textbf{Family} & \textbf{Statement} & \textbf{Deps} & \textbf{Order}\\
\midrule
\pidref{THM-LIMIT-UNIQUENESS} & Limits and Continuity & PASS & PASS & PASS \\
\pidref{THM-CHAIN-RULE} & Differential Calculus & PASS & PASS & PASS \\
\pidref{THM-FERMAT} & Optimization & PASS & PASS & PASS \\
\pidref{THM-FTC-I} & Integral Calculus & PASS & PASS & PASS \\
\pidref{THM-GEOMETRIC-SERIES} & Sequences and Series & PASS & PASS & PASS \\
\pidref{ALG-BISECTION-INVARIANT} & Root Finding & PASS & PASS & PASS \\
\pidref{THM-LAGRANGE-INTERPOLATION} & Numerical Approximation & PASS & PASS & PASS \\
\pidref{ERR-CENTERED-DIFFERENCE} & Numerical Approximation & PASS & PASS & PASS \\
\pidref{ERR-EULER-LOCAL} & Differential Equations & PASS & PASS & PASS \\
\pidref{THM-DIRECTIONAL-GRADIENT} & Multivariable Calculus & PASS & PASS & PASS \\
\bottomrule
\end{tabular}
\end{center}

\begin{landscape}
\section{Proof Registry and ID Index}
\scriptsize
\setlength{\tabcolsep}{2.6pt}
\begin{longtable}{@{}L{0.16\linewidth}L{0.15\linewidth}L{0.14\linewidth}L{0.25\linewidth}L{0.218\linewidth}C{0.04\linewidth}@{}}
\toprule
\textbf{Proof ID} & \textbf{Type} & \textbf{Family} & \textbf{Immediate dependencies} & \textbf{Preferred strategy} & \textbf{D}\\
\midrule
\endfirsthead
\toprule
\textbf{Proof ID} & \textbf{Type} & \textbf{Family} & \textbf{Immediate dependencies} & \textbf{Preferred strategy} & \textbf{D}\\
\midrule
\endhead
\pidraw{AX-CLASSICAL-LOGIC} & FOUNDATIONAL-ASSUMPTION & Foundations & - & Boundary & 0 \\
\pidraw{AX-COMPLEX-EXP} & FOUNDATIONAL-ASSUMPTION & Foundations & - & Boundary & 6 \\
\pidraw{AX-ELEMENTARY-SEQUENCE-LIMITS} & FOUNDATIONAL-ASSUMPTION & Foundations & - & Boundary & 2 \\
\pidraw{AX-EVT} & FOUNDATIONAL-ASSUMPTION & Foundations & - & Boundary & 3 \\
\pidraw{AX-LINEAR-ALGEBRA-BASICS} & FOUNDATIONAL-ASSUMPTION & Foundations & - & Boundary & 5 \\
\pidraw{AX-POLY-FACTOR} & FOUNDATIONAL-ASSUMPTION & Foundations & - & Boundary & 1 \\
\pidraw{AX-PROBABILITY} & FOUNDATIONAL-ASSUMPTION & Foundations & - & Boundary & 7 \\
\pidraw{AX-REAL-COMPLETENESS} & FOUNDATIONAL-ASSUMPTION & Foundations & - & Boundary & 1 \\
\pidraw{AX-REAL-FIELD} & FOUNDATIONAL-ASSUMPTION & Foundations & - & Boundary & 0 \\
\pidraw{AX-RIEMANN-BASIC} & FOUNDATIONAL-ASSUMPTION & Foundations & - & Boundary & 4 \\
\pidraw{DEF-BIG-O} & DEFINITION & Foundations & - & N/A & 2 \\
\pidraw{DEF-COMPOSITION} & DEFINITION & Foundations & - & N/A & 1 \\
\pidraw{DEF-CONTINUITY} & DEFINITION & Foundations & - & N/A & 1 \\
\pidraw{DEF-DERIVATIVE} & DEFINITION & Foundations & - & N/A & 2 \\
\pidraw{DEF-FUNCTION} & DEFINITION & Foundations & - & N/A & 1 \\
\pidraw{DEF-INVERSE-FUNCTION} & DEFINITION & Foundations & - & N/A & 1 \\
\pidraw{DEF-LIMIT-FUNCTION} & DEFINITION & Foundations & - & N/A & 1 \\
\pidraw{DEF-RIEMANN-INTEGRAL} & DEFINITION & Foundations & - & N/A & 4 \\
\pidraw{DEF-SERIES} & DEFINITION & Foundations & - & N/A & 5 \\
\pidraw{NOP-MODELING} & NOT-A-PROOF-OBJECT & Foundations & - & N/A & 0 \\
\pidraw{THM-NESTED-INTERVAL} & THEOREM & Limits and Continuity & AX-REAL-COMPLETENESS, AX-REAL-FIELD & Existence and uniqueness proof from completeness. & 3 \\
\pidraw{THM-IVT} & THEOREM & Limits and Continuity & DEF-CONTINUITY, AX-REAL-COMPLETENESS, AX-REAL-FIELD & Existence proof using completeness, followed by contradiction. & 3 \\
\pidraw{LEM-LIMIT-BOUNDED} & LEMMA & Limits and Continuity & DEF-LIMIT-FUNCTION, AX-REAL-FIELD & Direct proof from the definition of limit. & 1 \\
\pidraw{LEM-RECIPROCAL-BOUND} & LEMMA & Limits and Continuity & DEF-LIMIT-FUNCTION, AX-REAL-FIELD & Direct proof using the reverse triangle inequality. & 1 \\
\pidraw{THM-LIMIT-SUM} & THEOREM & Limits and Continuity & DEF-LIMIT-FUNCTION & Direct epsilon-delta proof. & 1 \\
\pidraw{THM-LIMIT-UNIQUENESS} & THEOREM & Limits and Continuity & DEF-LIMIT-FUNCTION, AX-REAL-FIELD & Proof by contradiction combined with a direct epsilon-delta estimate. & 1 \\
\pidraw{THM-LIMIT-PRODUCT} & THEOREM & Limits and Continuity & DEF-LIMIT-FUNCTION, LEM-LIMIT-BOUNDED, THM-LIMIT-SUM & Direct epsilon-delta proof with a local boundedness lemma. & 2 \\
\pidraw{THM-LIMIT-QUOTIENT} & THEOREM & Limits and Continuity & LEM-RECIPROCAL-BOUND, THM-LIMIT-SUM, THM-LIMIT-PRODUCT & Direct epsilon-delta proof using a denominator bound. & 2 \\
\pidraw{THM-DIFFERENTIABLE-CONTINUOUS} & THEOREM & Differential Calculus & DEF-DERIVATIVE, DEF-CONTINUITY, THM-LIMIT-PRODUCT & Direct limit proof. & 2 \\
\pidraw{THM-CHAIN-RULE} & THEOREM & Differential Calculus & DEF-DERIVATIVE, THM-DIFFERENTIABLE-CONTINUOUS, THM-LIMIT-PRODUCT & Direct limit proof using an auxiliary difference-quotient function. & 2 \\
\pidraw{THM-PRODUCT-RULE} & THEOREM & Differential Calculus & DEF-DERIVATIVE, THM-DIFFERENTIABLE-CONTINUOUS, THM-LIMIT-SUM, THM-LIMIT-PRODUCT & Direct proof from the derivative definition. & 2 \\
\pidraw{COR-IMPLICIT-DERIVATIVE} & COROLLARY & Differential Calculus & THM-CHAIN-RULE, DEF-DIFFERENTIABLE-RN & Algebraic derivation from the multivariable chain rule. & 3 \\
\pidraw{THM-QUOTIENT-RULE} & THEOREM & Differential Calculus & DEF-DERIVATIVE, THM-DIFFERENTIABLE-CONTINUOUS, THM-LIMIT-QUOTIENT, THM-PRODUCT-RULE & Direct derivative calculation using a controlled nonzero denominator. & 2 \\
\pidraw{DEF-CRITICAL-POINT} & DEFINITION & Optimization & - & N/A & 1 \\
\pidraw{THM-FERMAT} & THEOREM & Optimization & DEF-DERIVATIVE, AX-REAL-FIELD & Direct proof by one-sided sign comparison. & 3 \\
\pidraw{THM-ROLLE} & THEOREM & Optimization & AX-EVT, THM-FERMAT & Existence proof split into constant and nonconstant cases. & 3 \\
\pidraw{THM-MVT} & THEOREM & Optimization & THM-ROLLE & Direct reduction to Rolle's theorem. & 3 \\
\pidraw{COR-MONOTONICITY} & COROLLARY & Optimization & THM-MVT & Direct proof using the Mean Value Theorem. & 3 \\
\pidraw{THM-SECOND-DERIVATIVE-TEST} & THEOREM & Optimization & THM-MVT, COR-MONOTONICITY & Direct proof using continuity and two applications of derivative-sign reasoning. & 4 \\
\pidraw{DEF-RIEMANN-SUM} & DEFINITION & Integral Calculus & - & N/A & 1 \\
\pidraw{THM-FTC-I} & THEOREM & Integral Calculus & DEF-RIEMANN-INTEGRAL, DEF-CONTINUITY, AX-RIEMANN-BASIC & Direct epsilon estimate on the difference quotient. & 5 \\
\pidraw{THM-FTC-II} & THEOREM & Integral Calculus & THM-FTC-I, THM-MVT & Direct proof by comparing two functions with the same derivative. & 5 \\
\pidraw{DER-SUBSTITUTION} & DERIVATION & Integral Calculus & THM-CHAIN-RULE, THM-FTC-I, THM-FTC-II & Algebraic derivation from the chain rule and the Fundamental Theorem. & 5 \\
\pidraw{DER-INTEGRATION-BY-PARTS} & DERIVATION & Integral Calculus & THM-PRODUCT-RULE, THM-FTC-II & Algebraic derivation from the product rule. & 5 \\
\pidraw{THM-GEOMETRIC-SERIES} & THEOREM & Sequences and Series & DEF-SERIES, AX-REAL-FIELD, AX-ELEMENTARY-SEQUENCE-LIMITS & Direct proof through partial sums and a limit. & 5 \\
\pidraw{THM-SERIES-TERM-TEST} & THEOREM & Sequences and Series & DEF-SERIES, THM-LIMIT-SUM & Direct proof from consecutive partial sums. & 5 \\
\pidraw{THM-TAYLOR-LAGRANGE} & THEOREM & Sequences and Series & THM-ROLLE & Existence proof by constructing an auxiliary function and applying Rolle's theorem repeatedly. & 6 \\
\pidraw{DER-PARAMETRIC-DERIVATIVE} & DERIVATION & Parametric and Vector Calculus & THM-CHAIN-RULE & Algebraic derivation from the chain rule. & 6 \\
\pidraw{ADV-LAGRANGE} & SOURCE-BOUND FOUNDATION & Multivariable Calculus & - & Boundary & 10 \\
\pidraw{DEF-DIFFERENTIABLE-RN} & DEFINITION & Multivariable Calculus & - & N/A & 7 \\
\pidraw{DEF-DIRECTIONAL-DERIVATIVE} & DEFINITION & Multivariable Calculus & - & N/A & 7 \\
\pidraw{DEF-GRADIENT} & DEFINITION & Multivariable Calculus & - & N/A & 1 \\
\pidraw{DEF-PARTIAL-DERIVATIVE} & DEFINITION & Multivariable Calculus & - & N/A & 1 \\
\pidraw{THM-CURVE-CHAIN} & THEOREM & Multivariable Calculus & DEF-DIFFERENTIABLE-RN & Direct proof by composing the linear approximation with the curve increment. & 7 \\
\pidraw{THM-DIRECTIONAL-GRADIENT} & THEOREM & Multivariable Calculus & DEF-DIFFERENTIABLE-RN, DEF-DIRECTIONAL-DERIVATIVE & Direct proof from the definition of differentiability. & 7 \\
\pidraw{ADV-DIVERGENCE} & SOURCE-BOUND FOUNDATION & Vector Analysis & - & Boundary & 10 \\
\pidraw{ADV-GREEN} & SOURCE-BOUND FOUNDATION & Vector Analysis & - & Boundary & 10 \\
\pidraw{ADV-STOKES} & SOURCE-BOUND FOUNDATION & Vector Analysis & - & Boundary & 10 \\
\pidraw{DEF-LINE-INTEGRAL} & DEFINITION & Vector Analysis & - & N/A & 1 \\
\pidraw{DEF-VECTOR-FIELD} & DEFINITION & Vector Analysis & - & N/A & 1 \\
\pidraw{THM-FUNDAMENTAL-LINE-INTEGRAL} & THEOREM & Vector Analysis & THM-CURVE-CHAIN, THM-FTC-II & Direct reduction of the line integral to a one-variable derivative. & 8 \\
\pidraw{COR-PATH-INDEPENDENCE} & COROLLARY & Vector Analysis & THM-FUNDAMENTAL-LINE-INTEGRAL & Direct proof from endpoint evaluation. & 8 \\
\pidraw{ADV-CUBIC-SPLINE} & SOURCE-BOUND FOUNDATION & Numerical Approximation & - & Boundary & 10 \\
\pidraw{DEF-INTERPOLATION} & DEFINITION & Numerical Approximation & - & N/A & 1 \\
\pidraw{NOP-SPECIAL-FUNCTIONS} & NOT-A-PROOF-OBJECT & Numerical Approximation & - & N/A & 6 \\
\pidraw{LEM-POLYNOMIAL-ROOT-COUNT} & LEMMA & Numerical Approximation & AX-POLY-FACTOR & Mathematical induction on the degree. & 4 \\
\pidraw{THM-INTERPOLATION-UNIQUENESS} & THEOREM & Numerical Approximation & LEM-POLYNOMIAL-ROOT-COUNT & Proof by contradiction through a difference polynomial. & 4 \\
\pidraw{THM-LAGRANGE-INTERPOLATION} & THEOREM & Numerical Approximation & THM-INTERPOLATION-UNIQUENESS & Constructive existence proof followed by uniqueness. & 4 \\
\pidraw{ERR-CENTERED-DIFFERENCE} & ERROR-THEOREM & Numerical Approximation & THM-TAYLOR-LAGRANGE, DEF-BIG-O & Direct proof by symmetric Taylor expansions. & 5 \\
\pidraw{ERR-FORWARD-DIFFERENCE} & ERROR-THEOREM & Numerical Approximation & THM-TAYLOR-LAGRANGE, DEF-BIG-O & Direct Taylor expansion with remainder. & 5 \\
\pidraw{ADV-BRENT-ROOT} & SOURCE-BOUND FOUNDATION & Root Finding & - & Boundary & 8 \\
\pidraw{ADV-MULTIROOT} & SOURCE-BOUND FOUNDATION & Root Finding & - & Boundary & 10 \\
\pidraw{DER-NEWTON-ITERATION} & DERIVATION & Root Finding & DEF-DERIVATIVE & Algebraic derivation from the tangent-line equation. & 4 \\
\pidraw{ALG-SIGN-BRACKET-INVARIANT} & ALGORITHM-INVARIANT & Root Finding & THM-IVT & Direct proof from the Intermediate Value Theorem. & 4 \\
\pidraw{ALG-BISECTION-INVARIANT} & ALGORITHM-INVARIANT & Root Finding & ALG-SIGN-BRACKET-INVARIANT & Direct proof by cases. & 4 \\
\pidraw{CONV-BISECTION} & CONVERGENCE-THEOREM & Root Finding & ALG-BISECTION-INVARIANT, THM-NESTED-INTERVAL, DEF-CONTINUITY & Invariant plus limit argument. & 5 \\
\pidraw{CONV-NEWTON-LOCAL} & CONVERGENCE-THEOREM & Root Finding & DER-NEWTON-ITERATION, THM-TAYLOR-LAGRANGE & Direct Taylor remainder estimate. & 7 \\
\pidraw{ADV-BFGS} & SOURCE-BOUND FOUNDATION & Optimization Algorithms & - & Boundary & 10 \\
\pidraw{ADV-BRENT-MIN} & SOURCE-BOUND FOUNDATION & Optimization Algorithms & - & Boundary & 10 \\
\pidraw{ADV-NELDER-MEAD} & SOURCE-BOUND FOUNDATION & Optimization Algorithms & - & Boundary & 10 \\
\pidraw{ADV-SIMAN} & SOURCE-BOUND FOUNDATION & Optimization Algorithms & - & Boundary & 10 \\
\pidraw{CONV-GOLDEN-SECTION} & CONVERGENCE-THEOREM & Optimization Algorithms & AX-REAL-FIELD, AX-ELEMENTARY-SEQUENCE-LIMITS & Algebraic derivation followed by a geometric-sequence limit. & 5 \\
\pidraw{ADV-GAUSSIAN-QUADRATURE} & SOURCE-BOUND FOUNDATION & Numerical Integration & - & Boundary & 10 \\
\pidraw{ADV-LEBEDEV} & SOURCE-BOUND FOUNDATION & Numerical Integration & - & Boundary & 9 \\
\pidraw{DER-RICHARDSON} & DERIVATION & Numerical Integration & DEF-BIG-O & Algebraic derivation by cancellation of the leading error term. & 6 \\
\pidraw{THM-SIMPSON-EXACT-CUBIC} & THEOREM & Numerical Integration & THM-FTC-II & Direct proof by linearity and verification on a polynomial basis. & 5 \\
\pidraw{THM-TRAPEZOID-EXACT-LINEAR} & THEOREM & Numerical Integration & THM-FTC-II & Direct algebraic verification. & 5 \\
\pidraw{ADV-CG} & SOURCE-BOUND FOUNDATION & Linear Algebra & - & Boundary & 10 \\
\pidraw{ADV-SVD} & SOURCE-BOUND FOUNDATION & Linear Algebra & - & Boundary & 10 \\
\pidraw{DEF-EIGENPAIR} & DEFINITION & Linear Algebra & - & N/A & 1 \\
\pidraw{THM-LU-SOLVE-CORRECT} & THEOREM & Linear Algebra & AX-LINEAR-ALGEBRA-BASICS & Direct proof by substitution and invertibility. & 6 \\
\pidraw{THM-SYMMETRIC-EIGENVECTORS-ORTHOGONAL} & THEOREM & Linear Algebra & AX-LINEAR-ALGEBRA-BASICS & Direct algebraic proof. & 7 \\
\pidraw{THM-LEAST-SQUARES-NORMAL} & THEOREM & Linear Algebra & AX-LINEAR-ALGEBRA-BASICS, THM-FERMAT & Direct proof by orthogonality and quadratic expansion. & 7 \\
\pidraw{ADV-BVP} & SOURCE-BOUND FOUNDATION & Differential Equations & - & Boundary & 10 \\
\pidraw{ADV-RUNGE-KUTTA} & SOURCE-BOUND FOUNDATION & Differential Equations & - & Boundary & 10 \\
\pidraw{ADV-STIFF} & SOURCE-BOUND FOUNDATION & Differential Equations & - & Boundary & 10 \\
\pidraw{DEF-ODE-IVP} & DEFINITION & Differential Equations & - & N/A & 1 \\
\pidraw{ERR-EULER-LOCAL} & ERROR-THEOREM & Differential Equations & THM-TAYLOR-LAGRANGE, DEF-BIG-O & Direct Taylor expansion of the exact solution. & 7 \\
\pidraw{ADV-WAVELET} & SOURCE-BOUND FOUNDATION & Spectral Methods & - & Boundary & 10 \\
\pidraw{DEF-DFT} & DEFINITION & Spectral Methods & - & N/A & 1 \\
\pidraw{THM-DFT-EVEN-ODD} & THEOREM & Spectral Methods & AX-COMPLEX-EXP & Algebraic decomposition by parity. & 8 \\
\pidraw{THM-FFT-COMPLEXITY} & THEOREM & Spectral Methods & THM-DFT-EVEN-ODD, DEF-BIG-O & Mathematical induction on \$m=\textbackslash{}log\_2N\$. & 8 \\
\pidraw{NOP-DESCRIPTIVE-STATS} & NOT-A-PROOF-OBJECT & Probability and Monte Carlo & - & N/A & 7 \\
\pidraw{NOP-RNG} & NOT-A-PROOF-OBJECT & Probability and Monte Carlo & - & N/A & 8 \\
\pidraw{ERR-MC-VARIANCE} & ERROR-THEOREM & Probability and Monte Carlo & AX-PROBABILITY & Direct variance calculation. & 8 \\
\pidraw{THM-MC-UNBIASED} & THEOREM & Probability and Monte Carlo & AX-PROBABILITY & Direct proof by linearity of expectation. & 8 \\
\pidraw{ERR-SECOND-DERIVATIVE-CENTERED} & ERROR-THEOREM & PDE Discretization & THM-TAYLOR-LAGRANGE & Direct proof by symmetric fourth-order Taylor expansions. & 9 \\
\pidraw{NOP-COMPUTATIONAL-GEOMETRY} & NOT-A-PROOF-OBJECT & Geometry & - & N/A & 10 \\
\bottomrule
\end{longtable}
\end{landscape}

\begin{landscape}
\section{Numerical-Method Proof Coverage}
\small
\setlength{\tabcolsep}{3pt}
\begin{longtable}{@{}C{0.06\linewidth}L{0.23\linewidth}L{0.27\linewidth}L{0.38\linewidth}@{}}
\toprule
\textbf{Method} & \textbf{Catalog problem} & \textbf{\nr{} section} & \textbf{Proof / scope objects}\\
\midrule
\endfirsthead
\toprule
\textbf{Method} & \textbf{Catalog problem} & \textbf{\nr{} section} & \textbf{Proof / scope objects}\\
\midrule
\endhead
M00 & Floating-point error, accuracy, stability & NR 1.1 — Error, Accuracy, and Stability & \pidref{DEF-BIG-O}, \pidref{THM-LIMIT-UNIQUENESS} \\
M01 & Numerical first derivative & NR 5.7 — Numerical Derivatives & \pidref{ERR-FORWARD-DIFFERENCE}, \pidref{ERR-CENTERED-DIFFERENCE} \\
M02 & Polynomial interpolation & NR 3.2 — Polynomial Interpolation and Extrapolation & \pidref{THM-INTERPOLATION-UNIQUENESS}, \pidref{THM-LAGRANGE-INTERPOLATION} \\
M03 & Cubic-spline interpolation & NR 3.3 — Cubic Spline Interpolation & \pidref{ADV-CUBIC-SPLINE} \\
M04 & Bisection root finding & NR 9.1 — Bracketing and Bisection & \pidref{ALG-BISECTION-INVARIANT}, \pidref{CONV-BISECTION} \\
M05 & False-position root finding & NR 9.2 — Secant Method, False Position Method, and Ridders' Method & \pidref{ALG-SIGN-BRACKET-INVARIANT} \\
M06 & Brent–Dekker root finding & NR 9.3 — Van Wijngaarden-Dekker-Brent Method & \pidref{ALG-SIGN-BRACKET-INVARIANT}, \pidref{ADV-BRENT-ROOT} \\
M07 & Newton root polishing & NR 9.4 — Newton-Raphson Method Using Derivative & \pidref{DER-NEWTON-ITERATION}, \pidref{CONV-NEWTON-LOCAL} \\
M08 & Polynomial roots & NR 9.5 — Roots of Polynomials & \pidref{LEM-POLYNOMIAL-ROOT-COUNT} \\
M09 & Nonlinear systems & NR 9.6–9.7 — Newton-Raphson and globally convergent nonlinear systems & \pidref{ADV-MULTIROOT} \\
M10 & Golden-section scalar minimization & NR 10.2 — Golden Section Search in One Dimension & \pidref{CONV-GOLDEN-SECTION} \\
M11 & Brent scalar minimization & NR 10.3 — Parabolic Interpolation and Brent's Method in One Dimension & \pidref{ADV-BRENT-MIN} \\
M12 & Derivative-free multidimensional minimization & NR 10.5 — Downhill Simplex Method in Multidimensions & \pidref{ADV-NELDER-MEAD}, \pidref{ADV-LAGRANGE} \\
M13 & Conjugate-gradient minimization & NR 10.8 — Conjugate Gradient Methods in Multidimensions & \pidref{ADV-CG} \\
M14 & Quasi-Newton/BFGS minimization & NR 10.9 — Quasi-Newton or Variable Metric Methods & \pidref{ADV-BFGS}, \pidref{ADV-LAGRANGE} \\
M15 & Classical and Romberg quadrature & NR 4.1–4.3 — Classical formulas, elementary algorithms, Romberg & \pidref{THM-TRAPEZOID-EXACT-LINEAR}, \pidref{THM-SIMPSON-EXACT-CUBIC}, \pidref{DER-RICHARDSON} \\
M16 & Adaptive and improper quadrature & NR 4.4, 4.7 — Improper Integrals; Adaptive Quadrature & \pidref{DER-SUBSTITUTION}, \pidref{THM-FTC-I}, \pidref{THM-FTC-II} \\
M17 & Gaussian/fixed quadrature & NR 4.6 — Gaussian Quadratures and Orthogonal Polynomials & \pidref{ADV-GAUSSIAN-QUADRATURE} \\
M18 & Multidimensional Monte Carlo integration & NR 4.8; 7.7–7.9 — Multidimensional and Monte Carlo integration & \pidref{THM-MC-UNBIASED}, \pidref{ERR-MC-VARIANCE} \\
M19 & Spherical-surface quadrature & NR 4.8 — Multidimensional Integrals & \pidref{ADV-LEBEDEV} \\
M20 & Dense linear systems / LU & NR 2.3 — LU Decomposition and Its Applications & \pidref{THM-LU-SOLVE-CORRECT} \\
M21 & SVD / linear least squares & NR 2.6; 15.4 — SVD; General Linear Least Squares & \pidref{THM-LEAST-SQUARES-NORMAL}, \pidref{ADV-SVD} \\
M22 & Eigensystems & NR 11.1–11.8 — Eigensystems & \pidref{THM-SYMMETRIC-EIGENVECTORS-ORTHOGONAL} \\
M23 & Adaptive Runge–Kutta ODE integration & NR 17.1–17.2 — Runge-Kutta; adaptive stepsize control & \pidref{ERR-EULER-LOCAL}, \pidref{ADV-RUNGE-KUTTA}, \pidref{DER-RICHARDSON} \\
M24 & Stiff / implicit ODE integration & NR 17.5–17.6 — Stiff Sets; multistep/predictor-corrector & \pidref{ADV-STIFF} \\
M25 & Two-point boundary-value problems & NR 18.1–18.6 — Shooting and relaxation methods & \pidref{ADV-BVP} \\
M26 & Fast Fourier transform & NR 12.1–12.7 — Fast Fourier Transform & \pidref{THM-DFT-EVEN-ODD}, \pidref{THM-FFT-COMPLEXITY} \\
M27 & Fourier/spectral applications & NR 13.1–13.9 — Fourier and Spectral Applications & \pidref{THM-DFT-EVEN-ODD}, \pidref{THM-FFT-COMPLEXITY} \\
M28 & Descriptive statistics & NR 14.1–14.8 — Statistical Description of Data & \pidref{NOP-DESCRIPTIVE-STATS}, \pidref{THM-MC-UNBIASED} \\
M29 & Linear/nonlinear model fitting & NR 15.1–15.7 — Modeling of Data & \pidref{THM-LEAST-SQUARES-NORMAL} \\
M30 & Random sampling and stochastic simulation & NR 7.1–7.4 — Random Numbers & \pidref{NOP-RNG} \\
M31 & Special functions & NR 6.1–6.14 — Special Functions & \pidref{NOP-SPECIAL-FUNCTIONS} \\
M32 & Chebyshev approximation / series & NR 5.3, 5.8–5.10 — Series and Chebyshev approximation & \pidref{THM-GEOMETRIC-SERIES}, \pidref{THM-TAYLOR-LAGRANGE} \\
M33 & Finite-difference PDE linear systems & NR 20.1–20.7 — Partial Differential Equations & \pidref{ERR-SECOND-DERIVATIVE-CENTERED} \\
M34 & Integral equations / inverse problems & NR 19.1–19.7 — Integral Equations and Inverse Theory & \pidref{THM-LEAST-SQUARES-NORMAL}, \pidref{ADV-SVD} \\
M35 & Computational geometry & NR 21.1–21.8 — Computational Geometry & \pidref{NOP-COMPUTATIONAL-GEOMETRY} \\
M36 & Discrete wavelet transform & NR 13.10 — Wavelet Transforms & \pidref{ADV-WAVELET} \\
M37 & Simulated annealing & NR 10.12 — Simulated Annealing Methods & \pidref{ADV-SIMAN} \\
\bottomrule
\end{longtable}
\end{landscape}

\begin{landscape}
\section{Application-Technique Coverage Matrix}
\scriptsize
\setlength{\tabcolsep}{3pt}
\begin{longtable}{@{}L{0.39\linewidth}L{0.18\linewidth}L{0.37\linewidth}@{}}
\toprule
\textbf{Catalog mathematical technique} & \textbf{Applications} & \textbf{Canonical proof / definition objects}\\
\midrule
\endfirsthead
\toprule
\textbf{Catalog mathematical technique} & \textbf{Applications} & \textbf{Canonical proof / definition objects}\\
\midrule
\endhead
Arc-length integrals & A029 & \pidref{DEF-RIEMANN-INTEGRAL}, \pidref{THM-FTC-I}, \pidref{THM-FTC-II}, \pidref{DER-PARAMETRIC-DERIVATIVE} \\
Area between economic curves & A034 & \pidref{DEF-RIEMANN-INTEGRAL}, \pidref{THM-FTC-I}, \pidref{THM-FTC-II} \\
Average-value/multiple integration & A035 & \pidref{DEF-RIEMANN-INTEGRAL}, \pidref{THM-FTC-I}, \pidref{THM-FTC-II} \\
Boundary-value ODEs & A078 & \pidref{DEF-ODE-IVP}, \pidref{ERR-EULER-LOCAL}, \pidref{ADV-RUNGE-KUTTA}, \pidref{ADV-BVP} \\
Business optimization and constraints & A025 & \pidref{THM-FERMAT}, \pidref{THM-MVT}, \pidref{THM-SECOND-DERIVATIVE-TEST}, \pidref{CONV-GOLDEN-SECTION} \\
Chain rule and related rates & A008 & \pidref{THM-CHAIN-RULE} \\
Conics, intersections and computational geometry & A080 & \pidref{THM-IVT}, \pidref{ALG-SIGN-BRACKET-INVARIANT}, \pidref{DER-NEWTON-ITERATION}, \pidref{NOP-COMPUTATIONAL-GEOMETRY} \\
Constrained multivariable optimization & A049 & \pidref{THM-FERMAT}, \pidref{THM-MVT}, \pidref{THM-SECOND-DERIVATIVE-TEST}, \pidref{CONV-GOLDEN-SECTION}, \pidref{ADV-LAGRANGE} \\
Cost/objective optimization & A024 & \pidref{THM-FERMAT}, \pidref{THM-MVT}, \pidref{THM-SECOND-DERIVATIVE-TEST}, \pidref{CONV-GOLDEN-SECTION} \\
Coupled/separable population ODEs & A054 & \pidref{DEF-ODE-IVP}, \pidref{ERR-EULER-LOCAL}, \pidref{ADV-RUNGE-KUTTA}, \pidref{ADV-STIFF} \\
Definite integration and accumulation & A027 & \pidref{DEF-RIEMANN-INTEGRAL}, \pidref{THM-FTC-I}, \pidref{THM-FTC-II} \\
Derivative/integral rate relations & A010 & \pidref{DEF-RIEMANN-INTEGRAL}, \pidref{THM-FTC-I}, \pidref{THM-FTC-II} \\
Difference quotients, derivatives and higher derivatives & A006 & \pidref{DEF-DERIVATIVE}, \pidref{THM-DIFFERENTIABLE-CONTINUOUS} \\
Differentials and propagated error & A014 & \pidref{THM-DIFFERENTIABLE-CONTINUOUS}, \pidref{THM-TAYLOR-LAGRANGE} \\
Differentials, sensitivity and rate models & A013 & \pidref{THM-DIFFERENTIABLE-CONTINUOUS}, \pidref{THM-TAYLOR-LAGRANGE}, \pidref{DEF-ODE-IVP}, \pidref{ERR-EULER-LOCAL}, \pidref{ADV-RUNGE-KUTTA} \\
Divergence, flux and field equations & A077 & \pidref{ADV-DIVERGENCE}, \pidref{ADV-GREEN}, \pidref{ADV-STOKES}, \pidref{ERR-SECOND-DERIVATIVE-CENTERED} \\
Dot/cross products and vector geometry & A045 & \pidref{AX-LINEAR-ALGEBRA-BASICS} \\
Eigenvalue/spectral analysis & A072 & \pidref{DEF-EIGENPAIR}, \pidref{THM-SYMMETRIC-EIGENVECTORS-ORTHOGONAL}, \pidref{DEF-DFT}, \pidref{THM-DFT-EVEN-ODD}, \pidref{THM-FFT-COMPLEXITY} \\
Engineering objective functions & A022 & \pidref{THM-FERMAT}, \pidref{THM-MVT}, \pidref{THM-SECOND-DERIVATIVE-TEST}, \pidref{CONV-GOLDEN-SECTION} \\
Exponential decay ODE & A052 & \pidref{DEF-ODE-IVP}, \pidref{ERR-EULER-LOCAL}, \pidref{ADV-RUNGE-KUTTA} \\
Exponential models, series and accumulation & A041 & \pidref{DEF-SERIES}, \pidref{THM-GEOMETRIC-SERIES}, \pidref{THM-SERIES-TERM-TEST}, \pidref{THM-TAYLOR-LAGRANGE}, \pidref{DEF-ODE-IVP}, \pidref{ERR-EULER-LOCAL}, \pidref{ADV-RUNGE-KUTTA} \\
Extrema and curvature constraints & A023 & \pidref{THM-FERMAT}, \pidref{THM-MVT}, \pidref{THM-SECOND-DERIVATIVE-TEST}, \pidref{CONV-GOLDEN-SECTION} \\
Extrema in one and several variables & A020 & \pidref{THM-FERMAT}, \pidref{THM-MVT}, \pidref{THM-SECOND-DERIVATIVE-TEST}, \pidref{CONV-GOLDEN-SECTION} \\
Field equations and spectral/PDE methods & A079 & \pidref{DEF-DFT}, \pidref{THM-DFT-EVEN-ODD}, \pidref{THM-FFT-COMPLEXITY}, \pidref{ERR-SECOND-DERIVATIVE-CENTERED} \\
First-order and coupled circuit ODEs & A058 & \pidref{DEF-ODE-IVP}, \pidref{ERR-EULER-LOCAL}, \pidref{ADV-RUNGE-KUTTA}, \pidref{ADV-STIFF} \\
First-order linear compartment ODEs & A055 & \pidref{DEF-ODE-IVP}, \pidref{ERR-EULER-LOCAL}, \pidref{ADV-RUNGE-KUTTA} \\
First-order thermal ODEs & A056 & \pidref{DEF-ODE-IVP}, \pidref{ERR-EULER-LOCAL}, \pidref{ADV-RUNGE-KUTTA} \\
Fourier/spectral decomposition & A070 & \pidref{DEF-DFT}, \pidref{THM-DFT-EVEN-ODD}, \pidref{THM-FFT-COMPLEXITY} \\
Function modeling from data & A002 & \pidref{DEF-FUNCTION}, \pidref{NOP-MODELING}, \pidref{DEF-ODE-IVP}, \pidref{ERR-EULER-LOCAL}, \pidref{ADV-RUNGE-KUTTA} \\
Function modeling; sequences & A003 & \pidref{DEF-FUNCTION}, \pidref{NOP-MODELING}, \pidref{DEF-SERIES}, \pidref{THM-GEOMETRIC-SERIES}, \pidref{THM-SERIES-TERM-TEST}, \pidref{THM-TAYLOR-LAGRANGE}, \pidref{DEF-ODE-IVP}, \pidref{ERR-EULER-LOCAL}, \pidref{ADV-RUNGE-KUTTA} \\
Functional dependence and geometry & A005 & \pidref{DEF-FUNCTION}, \pidref{NOP-MODELING} \\
Functions, graphs, linear models and rates & A001 & \pidref{DEF-FUNCTION}, \pidref{NOP-MODELING}, \pidref{DEF-ODE-IVP}, \pidref{ERR-EULER-LOCAL}, \pidref{ADV-RUNGE-KUTTA} \\
Geometric optimization & A021 & \pidref{THM-FERMAT}, \pidref{THM-MVT}, \pidref{THM-SECOND-DERIVATIVE-TEST}, \pidref{CONV-GOLDEN-SECTION} \\
Geometry, optimization and multiple integration & A081 & \pidref{THM-FERMAT}, \pidref{THM-MVT}, \pidref{THM-SECOND-DERIVATIVE-TEST}, \pidref{CONV-GOLDEN-SECTION}, \pidref{DEF-RIEMANN-INTEGRAL}, \pidref{THM-FTC-I}, \pidref{THM-FTC-II} \\
Gradient-driven ODE path & A063 & \pidref{DEF-PARTIAL-DERIVATIVE}, \pidref{DEF-GRADIENT}, \pidref{THM-DIRECTIONAL-GRADIENT}, \pidref{DEF-ODE-IVP}, \pidref{ERR-EULER-LOCAL}, \pidref{ADV-RUNGE-KUTTA} \\
Growth/logistic ODEs & A053 & \pidref{DEF-ODE-IVP}, \pidref{ERR-EULER-LOCAL}, \pidref{ADV-RUNGE-KUTTA} \\
Hyperbolic PDEs & A075 & \pidref{ERR-SECOND-DERIVATIVE-CENTERED} \\
Implicit curves, special functions and quadrature & A039 & \pidref{COR-IMPLICIT-DERIVATIVE}, \pidref{NOP-SPECIAL-FUNCTIONS} \\
Implicit differentiation and related rates & A009 & \pidref{COR-IMPLICIT-DERIVATIVE} \\
Initial-value ODEs with drag & A060 & \pidref{DEF-ODE-IVP}, \pidref{ERR-EULER-LOCAL}, \pidref{ADV-RUNGE-KUTTA} \\
Integration of probability models & A036 & \pidref{DEF-RIEMANN-INTEGRAL}, \pidref{THM-FTC-I}, \pidref{THM-FTC-II}, \pidref{DEF-ODE-IVP}, \pidref{ERR-EULER-LOCAL}, \pidref{ADV-RUNGE-KUTTA}, \pidref{THM-MC-UNBIASED}, \pidref{ERR-MC-VARIANCE} \\
Interpolation, gradients and surfaces & A037 & \pidref{DEF-PARTIAL-DERIVATIVE}, \pidref{DEF-GRADIENT}, \pidref{THM-DIRECTIONAL-GRADIENT}, \pidref{THM-INTERPOLATION-UNIQUENESS}, \pidref{THM-LAGRANGE-INTERPOLATION}, \pidref{ADV-CUBIC-SPLINE} \\
Inverse problems and regularization & A083 & \pidref{THM-LEAST-SQUARES-NORMAL}, \pidref{ADV-SVD} \\
Linear/vector systems & A046 & \pidref{THM-LU-SOLVE-CORRECT}, \pidref{THM-LEAST-SQUARES-NORMAL} \\
Logarithmic/exponential functions & A040 & \pidref{NOP-SPECIAL-FUNCTIONS} \\
Moment integrals & A033 & \pidref{DEF-RIEMANN-INTEGRAL}, \pidref{THM-FTC-I}, \pidref{THM-FTC-II} \\
Multiple integrals and vector fields & A050 & \pidref{DEF-RIEMANN-INTEGRAL}, \pidref{THM-FTC-I}, \pidref{THM-FTC-II}, \pidref{DEF-VECTOR-FIELD}, \pidref{DEF-LINE-INTEGRAL}, \pidref{THM-FUNDAMENTAL-LINE-INTEGRAL}, \pidref{COR-PATH-INDEPENDENCE}, \pidref{AX-RIEMANN-BASIC} \\
Multivariable geometric optimization & A026 & \pidref{THM-FERMAT}, \pidref{THM-MVT}, \pidref{THM-SECOND-DERIVATIVE-TEST}, \pidref{CONV-GOLDEN-SECTION} \\
Nonlinear equations from geometry & A017 & \pidref{NOP-MODELING} \\
Nonlinear model fitting & A066 & \pidref{DEF-ODE-IVP}, \pidref{ERR-EULER-LOCAL}, \pidref{ADV-RUNGE-KUTTA}, \pidref{THM-LEAST-SQUARES-NORMAL}, \pidref{NOP-DESCRIPTIVE-STATS} \\
One- and multiple-integral volume & A028 & \pidref{DEF-RIEMANN-INTEGRAL}, \pidref{THM-FTC-I}, \pidref{THM-FTC-II} \\
Parabolic PDEs and boundary/initial conditions & A074 & \pidref{ERR-SECOND-DERIVATIVE-CENTERED} \\
Parametric/vector ODE trajectory models & A061 & \pidref{DEF-ODE-IVP}, \pidref{ERR-EULER-LOCAL}, \pidref{ADV-RUNGE-KUTTA} \\
Partial derivatives and gradients & A047 & \pidref{DEF-DERIVATIVE}, \pidref{THM-DIFFERENTIABLE-CONTINUOUS}, \pidref{DEF-PARTIAL-DERIVATIVE}, \pidref{DEF-GRADIENT}, \pidref{THM-DIRECTIONAL-GRADIENT} \\
Partial derivatives and marginal analysis & A011 & \pidref{DEF-DERIVATIVE}, \pidref{THM-DIFFERENTIABLE-CONTINUOUS}, \pidref{DEF-PARTIAL-DERIVATIVE}, \pidref{DEF-GRADIENT}, \pidref{THM-DIRECTIONAL-GRADIENT} \\
Partial derivatives, gradients and constrained extrema & A048 & \pidref{DEF-DERIVATIVE}, \pidref{THM-DIFFERENTIABLE-CONTINUOUS}, \pidref{THM-FERMAT}, \pidref{THM-MVT}, \pidref{THM-SECOND-DERIVATIVE-TEST}, \pidref{CONV-GOLDEN-SECTION}, \pidref{ADV-LAGRANGE}, \pidref{DEF-PARTIAL-DERIVATIVE}, \pidref{DEF-GRADIENT}, \pidref{THM-DIRECTIONAL-GRADIENT} \\
Periodic function analysis & A071 & \pidref{DEF-DFT}, \pidref{THM-DFT-EVEN-ODD}, \pidref{THM-FFT-COMPLEXITY} \\
Polynomial zeros & A018 & \pidref{THM-IVT}, \pidref{ALG-SIGN-BRACKET-INVARIANT}, \pidref{DER-NEWTON-ITERATION}, \pidref{LEM-POLYNOMIAL-ROOT-COUNT} \\
Position, velocity and acceleration & A007 & \pidref{DEF-DERIVATIVE}, \pidref{ERR-EULER-LOCAL} \\
Pressure-force integration & A032 & \pidref{DEF-RIEMANN-INTEGRAL}, \pidref{THM-FTC-I}, \pidref{THM-FTC-II} \\
Random variates and Monte Carlo & A069 & \pidref{THM-MC-UNBIASED}, \pidref{ERR-MC-VARIANCE} \\
Rates of change and exponential models & A012 & \pidref{DEF-DERIVATIVE}, \pidref{THM-DIFFERENTIABLE-CONTINUOUS}, \pidref{DEF-ODE-IVP}, \pidref{ERR-EULER-LOCAL}, \pidref{ADV-RUNGE-KUTTA} \\
Recursive geometry and spherical surfaces & A082 & \pidref{NOP-COMPUTATIONAL-GEOMETRY} \\
Regression, interpolation, and statistics & A068 & \pidref{THM-INTERPOLATION-UNIQUENESS}, \pidref{THM-LAGRANGE-INTERPOLATION}, \pidref{ADV-CUBIC-SPLINE}, \pidref{THM-LEAST-SQUARES-NORMAL}, \pidref{NOP-DESCRIPTIVE-STATS} \\
Regression, trend estimation and uncertainty & A064 & \pidref{THM-LEAST-SQUARES-NORMAL}, \pidref{NOP-DESCRIPTIVE-STATS} \\
Scalar and constrained optimization & A019 & \pidref{THM-FERMAT}, \pidref{THM-MVT}, \pidref{THM-SECOND-DERIVATIVE-TEST}, \pidref{CONV-GOLDEN-SECTION}, \pidref{ADV-LAGRANGE} \\
Scalar/vector fields and gradients & A051 & \pidref{DEF-VECTOR-FIELD}, \pidref{DEF-LINE-INTEGRAL}, \pidref{THM-FUNDAMENTAL-LINE-INTEGRAL}, \pidref{COR-PATH-INDEPENDENCE}, \pidref{DEF-PARTIAL-DERIVATIVE}, \pidref{DEF-GRADIENT}, \pidref{THM-DIRECTIONAL-GRADIENT} \\
Second-order ODEs; oscillation & A059 & \pidref{DEF-ODE-IVP}, \pidref{ERR-EULER-LOCAL}, \pidref{ADV-RUNGE-KUTTA} \\
Second-order gravitational ODEs & A062 & \pidref{DEF-ODE-IVP}, \pidref{ERR-EULER-LOCAL}, \pidref{ADV-RUNGE-KUTTA} \\
Separable ODEs and fluid models & A057 & \pidref{DEF-ODE-IVP}, \pidref{ERR-EULER-LOCAL}, \pidref{ADV-RUNGE-KUTTA} \\
Sequences, series and convergence & A042 & \pidref{DEF-SERIES}, \pidref{THM-GEOMETRIC-SERIES}, \pidref{THM-SERIES-TERM-TEST}, \pidref{THM-TAYLOR-LAGRANGE} \\
Solve nonlinear scalar equations & A016 & \pidref{THM-IVT}, \pidref{ALG-SIGN-BRACKET-INVARIANT}, \pidref{DER-NEWTON-ITERATION} \\
Special functions and nonlinear models & A043 & \pidref{DEF-ODE-IVP}, \pidref{ERR-EULER-LOCAL}, \pidref{ADV-RUNGE-KUTTA}, \pidref{NOP-SPECIAL-FUNCTIONS} \\
Statistical description and fitted models & A065 & \pidref{DEF-ODE-IVP}, \pidref{ERR-EULER-LOCAL}, \pidref{ADV-RUNGE-KUTTA}, \pidref{THM-LEAST-SQUARES-NORMAL}, \pidref{NOP-DESCRIPTIVE-STATS} \\
Statistics and fitted rates & A067 & \pidref{THM-LEAST-SQUARES-NORMAL}, \pidref{NOP-DESCRIPTIVE-STATS} \\
Surface-area integrals & A030 & \pidref{DEF-RIEMANN-INTEGRAL}, \pidref{THM-FTC-I}, \pidref{THM-FTC-II}, \pidref{DER-PARAMETRIC-DERIVATIVE} \\
Transcendental functions and integration & A038 & \pidref{DEF-RIEMANN-INTEGRAL}, \pidref{THM-FTC-I}, \pidref{THM-FTC-II}, \pidref{DER-SUBSTITUTION} \\
Trigonometric functions and parametric motion & A004 & \pidref{DER-PARAMETRIC-DERIVATIVE} \\
Vector addition and components & A044 & \pidref{AX-LINEAR-ALGEBRA-BASICS} \\
Vector calculus leading to PDEs & A076 & \pidref{ERR-SECOND-DERIVATIVE-CENTERED} \\
Wavelet transforms & A073 & \pidref{ADV-WAVELET} \\
Work integrals and line integrals & A031 & \pidref{DEF-RIEMANN-INTEGRAL}, \pidref{THM-FTC-I}, \pidref{THM-FTC-II}, \pidref{DEF-VECTOR-FIELD}, \pidref{DEF-LINE-INTEGRAL}, \pidref{THM-FUNDAMENTAL-LINE-INTEGRAL}, \pidref{COR-PATH-INDEPENDENCE} \\
Zeros/intersections of economic functions & A015 & \pidref{THM-IVT}, \pidref{ALG-SIGN-BRACKET-INVARIANT}, \pidref{DER-NEWTON-ITERATION} \\
\bottomrule
\end{longtable}
\end{landscape}
\section{Modular Referenced Proofs}

\subsection{Limits and Continuity}

\subsubsection{Nested intervals with vanishing lengths have one common point}

\begin{theorem}[Nested intervals with vanishing lengths have one common point]\label{prf:THM-NESTED-INTERVAL}

Let $I_n=[a_n,b_n]$ be nonempty closed intervals with
$I_{n+1}\subseteq I_n$ and $b_n-a_n\to0$. Then
$\bigcap_{n=1}^\infty I_n$ consists of exactly one real number.

\end{theorem}

\meta{Proof ID}{\pidraw{THM-NESTED-INTERVAL}}
\meta{Object type}{THEOREM}
\meta{Immediate dependencies}{\pidref{AX-REAL-COMPLETENESS}, \pidref{AX-REAL-FIELD}}
\meta{Additional transitive dependencies}{\emph{None}}
\meta{Source traceability}{\textbf{Proof method:} Mathematical Proofs, limits/real-number proof methodology; completeness boundary; \textbf{Numerical source:} Numerical Recipes, Sec. 9.1}
\meta{Catalog method IDs}{M04}

\begin{proofstrategy}Existence and uniqueness proof from completeness.\end{proofstrategy}

\begin{proof}

The sequence $(a_n)$ is nondecreasing and bounded above by every $b_k$.
Let $c=\sup\{a_n:n\in\mathbb N\}$. For each fixed $n$, all later left
endpoints satisfy $a_k\le b_n$, so $c\le b_n$; also $a_n\le c$. Thus
$c\in I_n$ for every $n$, proving existence.

If $c$ and $d$ both belong to every $I_n$, then
$|c-d|\le b_n-a_n$ for every $n$. Since the right side tends to $0$,
$|c-d|=0$, hence $c=d$. Thus the intersection is a singleton.

\end{proof}

\begin{proofanalysis}Completeness yields a common point from the monotone left endpoints; shrinking interval length forces uniqueness.\end{proofanalysis}

\begin{computationalrelevance}This is the topological core of bisection convergence.\end{computationalrelevance}

\subsubsection{Intermediate Value Theorem}

\begin{theorem}[Intermediate Value Theorem]\label{prf:THM-IVT}

Let $f:[a,b]\to\mathbb{R}$ be continuous. If a number $N$ lies between
$f(a)$ and $f(b)$, then there exists $c\in[a,b]$ such that $f(c)=N$.

\end{theorem}

\meta{Proof ID}{\pidraw{THM-IVT}}
\meta{Object type}{THEOREM}
\meta{Immediate dependencies}{\pidref{DEF-CONTINUITY}, \pidref{AX-REAL-COMPLETENESS}, \pidref{AX-REAL-FIELD}}
\meta{Additional transitive dependencies}{\emph{None}}
\meta{Source traceability}{\textbf{Proof method:} Mathematical Proofs, Sec. 5.4 for existence-proof structure; Ch. 14 for continuity; \textbf{Mathematical source:} source calculus framework, Sec. 1.4; \textbf{Numerical source:} Numerical Recipes, Sec. 9.1}
\meta{Catalog method IDs}{M04,M05,M06}

\begin{proofstrategy}Existence proof using completeness, followed by contradiction.\end{proofstrategy}

\begin{proof}

It suffices to treat the case $f(a)\le N\le f(b)$; the reverse ordering is analogous.
If either endpoint already maps to $N$, choose that endpoint. Otherwise
$f(a)<N<f(b)$. Define
\[
S=\{x\in[a,b]: f(x)<N\}.
\]
The set is nonempty because $a\in S$ and is bounded above by $b$, so completeness
gives $c=\sup S$.

Suppose first that $f(c)<N$. By continuity at $c$, there is $\eta>0$ such that
$|x-c|<\eta$ implies $f(x)<N$. Choosing a point
$x\in(c,\min\{c+\eta,b\})$ gives $x\in S$ with $x>c$, contradicting that $c$
is an upper bound.

Suppose instead that $f(c)>N$. By continuity there is $\eta>0$ such that
$|x-c|<\eta$ implies $f(x)>N$. By the definition of supremum there is
$s\in S$ with $c-\eta<s\le c$, but then $f(s)>N$, contradicting $s\in S$.
Therefore neither strict inequality is possible, and $f(c)=N$.

\end{proof}

\begin{proofanalysis}Completeness creates the candidate point; continuity prevents the function from jumping past the target value at that point.\end{proofanalysis}

\begin{computationalrelevance}The IVT is the existence theorem behind sign-bracketing root solvers such as bisection and safeguarded Brent methods.\end{computationalrelevance}

\subsubsection{A convergent function is locally bounded}

\begin{lemma}[A convergent function is locally bounded]\label{prf:LEM-LIMIT-BOUNDED}

If $\lim_{x\to a}f(x)=L$, then there are $\delta>0$ and $B>0$ such that
$0<|x-a|<\delta$ implies $|f(x)|\le B$.

\end{lemma}

\meta{Proof ID}{\pidraw{LEM-LIMIT-BOUNDED}}
\meta{Object type}{LEMMA}
\meta{Immediate dependencies}{\pidref{DEF-LIMIT-FUNCTION}, \pidref{AX-REAL-FIELD}}
\meta{Additional transitive dependencies}{\emph{None}}
\meta{Source traceability}{\textbf{Proof method:} Mathematical Proofs, Sec. 14.4}

\begin{proofstrategy}Direct proof from the definition of limit.\end{proofstrategy}

\begin{proof}

Use the definition with $\varepsilon=1$. Then for some $\delta>0$,
$0<|x-a|<\delta$ implies $|f(x)-L|<1$. Hence
\[
|f(x)|\le |f(x)-L|+|L|<1+|L|.
\]
Thus $B=1+|L|$ works.

\end{proof}

\begin{proofanalysis}A finite limit constrains nearby function values to a finite interval.\end{proofanalysis}

\begin{computationalrelevance}The boundedness estimate is the key technical step in the product limit law and in many numerical error estimates.\end{computationalrelevance}

\subsubsection{A nonzero limit keeps the denominator away from zero}

\begin{lemma}[A nonzero limit keeps the denominator away from zero]\label{prf:LEM-RECIPROCAL-BOUND}

If $\lim_{x\to a}g(x)=M$ with $M\ne0$, then for all $x$ in some deleted neighborhood of $a$,
\[
|g(x)|>\frac{|M|}{2},\qquad \frac1{|g(x)|}<\frac2{|M|}.
\]

\end{lemma}

\meta{Proof ID}{\pidraw{LEM-RECIPROCAL-BOUND}}
\meta{Object type}{LEMMA}
\meta{Immediate dependencies}{\pidref{DEF-LIMIT-FUNCTION}, \pidref{AX-REAL-FIELD}}
\meta{Additional transitive dependencies}{\emph{None}}
\meta{Source traceability}{\textbf{Proof method:} Mathematical Proofs, Lemma 14.26}

\begin{proofstrategy}Direct proof using the reverse triangle inequality.\end{proofstrategy}

\begin{proof}

Take $\varepsilon=|M|/2$. Then for sufficiently small $0<|x-a|$,
$|g(x)-M|<|M|/2$. Since
\[
|M|\le |M-g(x)|+|g(x)|,
\]
we obtain
\[
|g(x)|>|M|-\frac{|M|}{2}=\frac{|M|}{2}.
\]
Taking reciprocals of positive quantities yields the stated reciprocal bound.

\end{proof}

\begin{proofanalysis}The point is not merely that $g(x)$ is nonzero; it is uniformly separated from zero near $a$.\end{proofanalysis}

\begin{computationalrelevance}Division, rational interpolation, and Newton-type expressions require a controlled nonzero denominator.\end{computationalrelevance}

\subsubsection{Sum law for limits}

\begin{theorem}[Sum law for limits]\label{prf:THM-LIMIT-SUM}

If $\lim_{x\to a}f(x)=L$ and $\lim_{x\to a}g(x)=M$, then
\[
\lim_{x\to a}(f(x)+g(x))=L+M.
\]

\end{theorem}

\meta{Proof ID}{\pidraw{THM-LIMIT-SUM}}
\meta{Object type}{THEOREM}
\meta{Immediate dependencies}{\pidref{DEF-LIMIT-FUNCTION}}
\meta{Additional transitive dependencies}{\emph{None}}
\meta{Source traceability}{\textbf{Proof method:} Mathematical Proofs, Theorem 14.23}

\begin{proofstrategy}Direct epsilon-delta proof.\end{proofstrategy}

\begin{proof}

Let $\varepsilon>0$. Choose $\delta_f>0$ so that
$|f(x)-L|<\varepsilon/2$ whenever $0<|x-a|<\delta_f$, and choose
$\delta_g>0$ similarly for $g$ and $M$. Set
$\delta=\min(\delta_f,\delta_g)$. Then
\[
|(f(x)+g(x))-(L+M)|
\le |f(x)-L|+|g(x)-M|<\varepsilon.
\]
Therefore the asserted limit holds.

\end{proof}

\begin{proofanalysis}The error budget is split into two halves and synchronized by the minimum of the two neighborhoods.\end{proofanalysis}

\begin{computationalrelevance}This law underlies algebraic manipulation of truncation errors, interpolation formulas, and discretizations.\end{computationalrelevance}

\subsubsection{Uniqueness of a function limit}

\begin{theorem}[Uniqueness of a function limit]\label{prf:THM-LIMIT-UNIQUENESS}

If $\lim_{x\to a}f(x)=L$ and $\lim_{x\to a}f(x)=M$, then $L=M$.

\end{theorem}

\meta{Proof ID}{\pidraw{THM-LIMIT-UNIQUENESS}}
\meta{Object type}{THEOREM}
\meta{Hypotheses}{Both limits exist in the $\varepsilon$--$\delta$ sense.}
\meta{Conclusion}{The two candidate limits are equal.}
\meta{Immediate dependencies}{\pidref{DEF-LIMIT-FUNCTION}, \pidref{AX-REAL-FIELD}}
\meta{Additional transitive dependencies}{\emph{None}}
\meta{Source traceability}{\textbf{Proof method:} Mathematical Proofs, Sec. 5.2 and Ch. 14; \textbf{Mathematical source:} source calculus framework, Sec. 1.3--1.4}

\begin{proofstrategy}Proof by contradiction combined with a direct epsilon-delta estimate.\end{proofstrategy}

\begin{proof}

Assume, to the contrary, that $L\ne M$. Put $\varepsilon=|L-M|/3>0$.
There are $\delta_L,\delta_M>0$ such that
\[
0<|x-a|<\delta_L\Rightarrow |f(x)-L|<\varepsilon,
\qquad
0<|x-a|<\delta_M\Rightarrow |f(x)-M|<\varepsilon.
\]
Choose $x$ in the deleted neighborhood with
$0<|x-a|<\min(\delta_L,\delta_M)$. Then the triangle inequality gives
\[
|L-M|\le |L-f(x)|+|f(x)-M|<2\varepsilon=\frac23|L-M|,
\]
a contradiction. Hence $L=M$.

\end{proof}

\begin{proofanalysis}The contradiction is created by forcing one function value to lie simultaneously in two disjoint sufficiently small neighborhoods of distinct numbers.\end{proofanalysis}

\begin{computationalrelevance}A numerical method can speak of ``the'' limiting value only because a limit, when it exists, is unique.\end{computationalrelevance}

\subsubsection{Product law for limits}

\begin{theorem}[Product law for limits]\label{prf:THM-LIMIT-PRODUCT}

If $\lim_{x\to a}f(x)=L$ and $\lim_{x\to a}g(x)=M$, then
\[
\lim_{x\to a}f(x)g(x)=LM.
\]

\end{theorem}

\meta{Proof ID}{\pidraw{THM-LIMIT-PRODUCT}}
\meta{Object type}{THEOREM}
\meta{Immediate dependencies}{\pidref{DEF-LIMIT-FUNCTION}, \pidref{LEM-LIMIT-BOUNDED}, \pidref{THM-LIMIT-SUM}}
\meta{Additional transitive dependencies}{\pidref{AX-REAL-FIELD}}
\meta{Source traceability}{\textbf{Proof method:} Mathematical Proofs, Theorem 14.25}

\begin{proofstrategy}Direct epsilon-delta proof with a local boundedness lemma.\end{proofstrategy}

\begin{proof}

By \pid{LEM-LIMIT-BOUNDED}, choose a deleted neighborhood on which
$|f(x)|\le B$ for some $B>0$. For $\varepsilon>0$, choose neighborhoods so that
\[
|g(x)-M|<\frac{\varepsilon}{2B},
\]
and, when $M\ne0$,
\[
|f(x)-L|<\frac{\varepsilon}{2|M|}.
\]
If $M=0$, the second requirement is unnecessary. On the common neighborhood,
\[
|f(x)g(x)-LM|
\le |f(x)|\,|g(x)-M|+|M|\,|f(x)-L|<\varepsilon.
\]
Thus $f(x)g(x)\to LM$.

\end{proof}

\begin{proofanalysis}Adding and subtracting $f(x)M$ separates the product error into two controllable terms.\end{proofanalysis}

\begin{computationalrelevance}Products of approximations occur throughout quadrature, differential equations, and error propagation.\end{computationalrelevance}

\subsubsection{Quotient law for limits}

\begin{theorem}[Quotient law for limits]\label{prf:THM-LIMIT-QUOTIENT}

If $\lim_{x\to a}f(x)=L$ and $\lim_{x\to a}g(x)=M\ne0$, then
\[
\lim_{x\to a}\frac{f(x)}{g(x)}=\frac{L}{M}.
\]

\end{theorem}

\meta{Proof ID}{\pidraw{THM-LIMIT-QUOTIENT}}
\meta{Object type}{THEOREM}
\meta{Immediate dependencies}{\pidref{LEM-RECIPROCAL-BOUND}, \pidref{THM-LIMIT-SUM}, \pidref{THM-LIMIT-PRODUCT}}
\meta{Additional transitive dependencies}{\pidref{AX-REAL-FIELD}, \pidref{DEF-LIMIT-FUNCTION}, \pidref{LEM-LIMIT-BOUNDED}}
\meta{Source traceability}{\textbf{Proof method:} Mathematical Proofs, Lemma 14.26 and Theorem 14.27}

\begin{proofstrategy}Direct epsilon-delta proof using a denominator bound.\end{proofstrategy}

\begin{proof}

By \pid{LEM-RECIPROCAL-BOUND}, in a sufficiently small deleted neighborhood
$1/|g(x)|<2/|M|$. Write
\[
\left|\frac{f(x)}{g(x)}-\frac{L}{M}\right|
=
\frac{|M(f(x)-L)+L(M-g(x))|}{|g(x)||M|}.
\]
Hence
\[
\left|\frac{f(x)}{g(x)}-\frac{L}{M}\right|
\le
\frac{2}{|M|}|f(x)-L|
+\frac{2|L|}{|M|^2}|g(x)-M|.
\]
For a given $\varepsilon>0$, make the first term less than $\varepsilon/2$
and, when $L\ne0$, the second term less than $\varepsilon/2$ by choosing the
common neighborhood sufficiently small. If $L=0$, only the first term remains.
The quotient therefore tends to $L/M$.

\end{proof}

\begin{proofanalysis}The denominator lemma converts division into a pair of ordinary limit estimates.\end{proofanalysis}

\begin{computationalrelevance}Rational formulas occur in Newton updates, relative error measures, and many special-function evaluations.\end{computationalrelevance}

\subsection{Differential Calculus}

\subsubsection{Differentiability implies continuity}

\begin{theorem}[Differentiability implies continuity]\label{prf:THM-DIFFERENTIABLE-CONTINUOUS}

If $f$ is differentiable at $a$, then $f$ is continuous at $a$.

\end{theorem}

\meta{Proof ID}{\pidraw{THM-DIFFERENTIABLE-CONTINUOUS}}
\meta{Object type}{THEOREM}
\meta{Immediate dependencies}{\pidref{DEF-DERIVATIVE}, \pidref{DEF-CONTINUITY}, \pidref{THM-LIMIT-PRODUCT}}
\meta{Additional transitive dependencies}{\pidref{AX-REAL-FIELD}, \pidref{DEF-LIMIT-FUNCTION}, \pidref{LEM-LIMIT-BOUNDED}, \pidref{THM-LIMIT-SUM}}
\meta{Source traceability}{\textbf{Proof method:} Mathematical Proofs, Theorem 14.38; \textbf{Mathematical source:} source calculus framework, Secs. 2.1 and 3}

\begin{proofstrategy}Direct limit proof.\end{proofstrategy}

\begin{proof}

For $x\ne a$,
\[
f(x)-f(a)=
\frac{f(x)-f(a)}{x-a}(x-a).
\]
As $x\to a$, the first factor tends to $f'(a)$ by differentiability and the
second tends to $0$. By \pid{THM-LIMIT-PRODUCT},
$f(x)-f(a)\to0$. Hence $f(x)\to f(a)$, which is continuity at $a$.

\end{proof}

\begin{proofanalysis}The increment in the function factors into a bounded slope-like term and a vanishing horizontal increment.\end{proofanalysis}

\begin{computationalrelevance}Many optimization and root-finding theorems require continuity; differentiability supplies it automatically.\end{computationalrelevance}

\subsubsection{Chain rule}

\begin{theorem}[Chain rule]\label{prf:THM-CHAIN-RULE}

If $g$ is differentiable at $a$ and $f$ is differentiable at $g(a)$, then $f\circ g$ is differentiable at $a$ and
\[
(f\circ g)'(a)=f'(g(a))\,g'(a).
\]

\end{theorem}

\meta{Proof ID}{\pidraw{THM-CHAIN-RULE}}
\meta{Object type}{THEOREM}
\meta{Immediate dependencies}{\pidref{DEF-DERIVATIVE}, \pidref{THM-DIFFERENTIABLE-CONTINUOUS}, \pidref{THM-LIMIT-PRODUCT}}
\meta{Additional transitive dependencies}{\pidref{AX-REAL-FIELD}, \pidref{DEF-CONTINUITY}, \pidref{DEF-LIMIT-FUNCTION}, \pidref{LEM-LIMIT-BOUNDED}, \pidref{THM-LIMIT-SUM}}
\meta{Source traceability}{\textbf{Proof method:} Mathematical Proofs, direct proof and Sec. 14.6; \textbf{Mathematical source:} source calculus framework, Sec. 2.4}

\begin{proofstrategy}Direct limit proof using an auxiliary difference-quotient function.\end{proofstrategy}

\begin{proof}

Let $b=g(a)$. Define
\[
\phi(y)=
\begin{cases}
\dfrac{f(y)-f(b)}{y-b},&y\ne b,\\[4pt]
f'(b),&y=b.
\end{cases}
\]
Differentiability of $f$ at $b$ means $\phi(y)\to f'(b)$ as $y\to b$.
Since $g$ is differentiable at $a$, it is continuous there, so $g(x)\to b$ and
therefore $\phi(g(x))\to f'(b)$. For $x\ne a$,
\[
\frac{f(g(x))-f(g(a))}{x-a}
=
\phi(g(x))\frac{g(x)-g(a)}{x-a}.
\]
Taking limits gives $f'(b)g'(a)$ by the product limit law.

\end{proof}

\begin{proofanalysis}The auxiliary function removes the usual difficulty caused by points where $g(x)=g(a)$.\end{proofanalysis}

\begin{computationalrelevance}The chain rule is central to implicit differentiation, sensitivity propagation, nonlinear optimization, and Newton-type methods.\end{computationalrelevance}

\subsubsection{Product rule}

\begin{theorem}[Product rule]\label{prf:THM-PRODUCT-RULE}

If $f$ and $g$ are differentiable at $a$, then $fg$ is differentiable at $a$ and
\[
(fg)'(a)=f'(a)g(a)+f(a)g'(a).
\]

\end{theorem}

\meta{Proof ID}{\pidraw{THM-PRODUCT-RULE}}
\meta{Object type}{THEOREM}
\meta{Immediate dependencies}{\pidref{DEF-DERIVATIVE}, \pidref{THM-DIFFERENTIABLE-CONTINUOUS}, \pidref{THM-LIMIT-SUM}, \pidref{THM-LIMIT-PRODUCT}}
\meta{Additional transitive dependencies}{\pidref{AX-REAL-FIELD}, \pidref{DEF-CONTINUITY}, \pidref{DEF-LIMIT-FUNCTION}, \pidref{LEM-LIMIT-BOUNDED}}
\meta{Source traceability}{\textbf{Proof method:} Mathematical Proofs, direct-proof methodology and Sec. 14.6; \textbf{Mathematical source:} source calculus framework, Sec. 2.3}

\begin{proofstrategy}Direct proof from the derivative definition.\end{proofstrategy}

\begin{proof}

For $x\ne a$, add and subtract $f(x)g(a)$:
\[
\frac{f(x)g(x)-f(a)g(a)}{x-a}
=
f(x)\frac{g(x)-g(a)}{x-a}
+
g(a)\frac{f(x)-f(a)}{x-a}.
\]
Differentiability of $f$ implies $f(x)\to f(a)$ by
\pid{THM-DIFFERENTIABLE-CONTINUOUS}. Taking limits and applying the sum and
product limit laws gives
\[
(fg)'(a)=f(a)g'(a)+g(a)f'(a).
\]

\end{proof}

\begin{proofanalysis}The inserted intermediate product turns one difference of products into two derivative quotients.\end{proofanalysis}

\begin{computationalrelevance}The rule supports integration by parts, Jacobian calculations, sensitivity analysis, and many model derivatives.\end{computationalrelevance}

\subsubsection{Implicit derivative formula}

\begin{corollary}[Implicit derivative formula]\label{prf:COR-IMPLICIT-DERIVATIVE}

Suppose $F(x,y(x))=0$, $F$ is differentiable, $y$ is differentiable, and
$F_y(x,y(x))\ne0$. Then
\[
y'(x)=-\frac{F_x(x,y(x))}{F_y(x,y(x))}.
\]

\end{corollary}

\meta{Proof ID}{\pidraw{COR-IMPLICIT-DERIVATIVE}}
\meta{Object type}{COROLLARY}
\meta{Immediate dependencies}{\pidref{THM-CHAIN-RULE}, \pidref{DEF-DIFFERENTIABLE-RN}}
\meta{Additional transitive dependencies}{\pidref{AX-REAL-FIELD}, \pidref{DEF-CONTINUITY}, \pidref{DEF-DERIVATIVE}, \pidref{DEF-LIMIT-FUNCTION}, \pidref{LEM-LIMIT-BOUNDED}, \pidref{THM-LIMIT-SUM}, \pidref{THM-LIMIT-PRODUCT}, \pidref{THM-DIFFERENTIABLE-CONTINUOUS}}
\meta{Source traceability}{\textbf{Mathematical source:} source calculus framework, Secs. 2.5 and 13.5}

\begin{proofstrategy}Algebraic derivation from the multivariable chain rule.\end{proofstrategy}

\begin{proof}

Differentiate the identity $F(x,y(x))=0$ with respect to $x$. The chain rule for
a differentiable function of two variables along the differentiable curve
$x\mapsto(x,y(x))$ gives
\[
F_x(x,y(x))+F_y(x,y(x))y'(x)=0.
\]
Because $F_y(x,y(x))\ne0$, solving for $y'(x)$ yields
\[
y'(x)=-\frac{F_x(x,y(x))}{F_y(x,y(x))}.
\]

\end{proof}

\begin{proofanalysis}The nonzero partial derivative is exactly the algebraic condition needed to isolate the unknown rate.\end{proofanalysis}

\begin{computationalrelevance}Implicit equations arise in related-rates problems, calibration equations, constrained geometry, and nonlinear systems.\end{computationalrelevance}

\subsubsection{Quotient rule}

\begin{theorem}[Quotient rule]\label{prf:THM-QUOTIENT-RULE}

If $f$ and $g$ are differentiable at $a$ and $g(a)\ne0$, then $f/g$ is differentiable at $a$ and
\[
\left(\frac{f}{g}\right)'(a)=
\frac{f'(a)g(a)-f(a)g'(a)}{g(a)^2}.
\]

\end{theorem}

\meta{Proof ID}{\pidraw{THM-QUOTIENT-RULE}}
\meta{Object type}{THEOREM}
\meta{Immediate dependencies}{\pidref{DEF-DERIVATIVE}, \pidref{THM-DIFFERENTIABLE-CONTINUOUS}, \pidref{THM-LIMIT-QUOTIENT}, \pidref{THM-PRODUCT-RULE}}
\meta{Additional transitive dependencies}{\pidref{AX-REAL-FIELD}, \pidref{DEF-CONTINUITY}, \pidref{DEF-LIMIT-FUNCTION}, \pidref{LEM-LIMIT-BOUNDED}, \pidref{LEM-RECIPROCAL-BOUND}, \pidref{THM-LIMIT-SUM}, \pidref{THM-LIMIT-PRODUCT}}
\meta{Source traceability}{\textbf{Proof method:} Mathematical Proofs, direct-proof methodology and Sec. 14.6; \textbf{Mathematical source:} source calculus framework, Sec. 2.3}

\begin{proofstrategy}Direct derivative calculation using a controlled nonzero denominator.\end{proofstrategy}

\begin{proof}

By continuity of $g$ at $a$ and $g(a)\ne0$, $g(x)\ne0$ for all $x$ sufficiently
near $a$. For such $x\ne a$,
\[
\frac{\frac{f(x)}{g(x)}-\frac{f(a)}{g(a)}}{x-a}
=
\frac{g(a)\frac{f(x)-f(a)}{x-a}
      -f(a)\frac{g(x)-g(a)}{x-a}}
     {g(x)g(a)}.
\]
As $x\to a$, the numerator tends to
$g(a)f'(a)-f(a)g'(a)$ and the denominator tends to $g(a)^2\ne0$.
The quotient limit law therefore gives the stated derivative.

\end{proof}

\begin{proofanalysis}The algebra isolates the two derivative quotients while continuity guarantees the denominator remains valid.\end{proofanalysis}

\begin{computationalrelevance}Rational response models, relative sensitivities, and Newton updates frequently require quotient differentiation.\end{computationalrelevance}

\subsection{Optimization}

\subsubsection{Necessary derivative condition at an interior extremum}

\begin{theorem}[Necessary derivative condition at an interior extremum]\label{prf:THM-FERMAT}

If $f$ has a local maximum or local minimum at an interior point $c$ and $f'(c)$ exists, then $f'(c)=0$.

\end{theorem}

\meta{Proof ID}{\pidraw{THM-FERMAT}}
\meta{Object type}{THEOREM}
\meta{Immediate dependencies}{\pidref{DEF-DERIVATIVE}, \pidref{AX-REAL-FIELD}}
\meta{Additional transitive dependencies}{\emph{None}}
\meta{Source traceability}{\textbf{Proof method:} Mathematical Proofs, direct proof / proof by cases; \textbf{Mathematical source:} source calculus framework, Secs. 3.1 and 3.7}

\begin{proofstrategy}Direct proof by one-sided sign comparison.\end{proofstrategy}

\begin{proof}

Assume $f$ has a local minimum at $c$; the maximum case is analogous.
For all sufficiently small $h$,
$f(c+h)-f(c)\ge0$. If $h>0$, then
\[
\frac{f(c+h)-f(c)}{h}\ge0,
\]
whereas if $h<0$, then
\[
\frac{f(c+h)-f(c)}{h}\le0.
\]
If the two-sided derivative exists, the right and left limits must agree.
A real number that is both at least $0$ and at most $0$ is $0$.
Therefore $f'(c)=0$.

\end{proof}

\begin{proofanalysis}An interior extremum forces opposite signs on the one-sided difference quotients; differentiability forces them to meet.\end{proofanalysis}

\begin{computationalrelevance}This is the analytical basis for stationary-point searches and for reducing smooth optimization to root finding on derivatives.\end{computationalrelevance}

\subsubsection{Rolle's theorem}

\begin{theorem}[Rolle's theorem]\label{prf:THM-ROLLE}

If $f$ is continuous on $[a,b]$, differentiable on $(a,b)$, and $f(a)=f(b)$, then there exists $c\in(a,b)$ such that $f'(c)=0$.

\end{theorem}

\meta{Proof ID}{\pidraw{THM-ROLLE}}
\meta{Object type}{THEOREM}
\meta{Immediate dependencies}{\pidref{AX-EVT}, \pidref{THM-FERMAT}}
\meta{Additional transitive dependencies}{\pidref{AX-REAL-FIELD}, \pidref{DEF-DERIVATIVE}}
\meta{Source traceability}{\textbf{Proof method:} Mathematical Proofs, existence-proof structure and calculus-proof methodology; \textbf{Mathematical source:} source calculus framework, Sec. 3.2}

\begin{proofstrategy}Existence proof split into constant and nonconstant cases.\end{proofstrategy}

\begin{proof}

If $f$ is constant on $[a,b]$, every $c\in(a,b)$ satisfies $f'(c)=0$.
Otherwise, by the Extreme Value Theorem, $f$ attains a maximum and a minimum
on $[a,b]$. Because $f(a)=f(b)$ and the function is not constant, at least one
of those extreme values is attained at an interior point $c\in(a,b)$.
By \pid{THM-FERMAT}, $f'(c)=0$.

\end{proof}

\begin{proofanalysis}Compact-interval existence of an extremum is converted into a derivative-zero conclusion.\end{proofanalysis}

\begin{computationalrelevance}Rolle's theorem is the structural lemma behind the Mean Value Theorem and Taylor remainder arguments.\end{computationalrelevance}

\subsubsection{Mean Value Theorem}

\begin{theorem}[Mean Value Theorem]\label{prf:THM-MVT}

If $f$ is continuous on $[a,b]$ and differentiable on $(a,b)$, then there exists $c\in(a,b)$ such that
\[
f'(c)=\frac{f(b)-f(a)}{b-a}.
\]

\end{theorem}

\meta{Proof ID}{\pidraw{THM-MVT}}
\meta{Object type}{THEOREM}
\meta{Immediate dependencies}{\pidref{THM-ROLLE}}
\meta{Additional transitive dependencies}{\pidref{AX-EVT}, \pidref{AX-REAL-FIELD}, \pidref{DEF-DERIVATIVE}, \pidref{THM-FERMAT}}
\meta{Source traceability}{\textbf{Proof method:} Mathematical Proofs, direct proof and existence proof; \textbf{Mathematical source:} source calculus framework, Sec. 3.2}

\begin{proofstrategy}Direct reduction to Rolle's theorem.\end{proofstrategy}

\begin{proof}

Define the secant-line correction
\[
g(x)=f(x)-\frac{f(b)-f(a)}{b-a}(x-a).
\]
Then $g$ is continuous on $[a,b]$, differentiable on $(a,b)$, and
$g(a)=f(a)=g(b)$. By \pid{THM-ROLLE}, there is $c\in(a,b)$ with $g'(c)=0$.
Therefore
\[
0=f'(c)-\frac{f(b)-f(a)}{b-a},
\]
which is the asserted formula.

\end{proof}

\begin{proofanalysis}Subtracting the secant line converts an arbitrary endpoint difference into equal endpoint values.\end{proofanalysis}

\begin{computationalrelevance}The MVT supports monotonicity tests, error bounds, Newton convergence estimates, and first-order sensitivity reasoning.\end{computationalrelevance}

\subsubsection{Derivative sign implies monotonicity}

\begin{corollary}[Derivative sign implies monotonicity]\label{prf:COR-MONOTONICITY}

If $f$ is differentiable on an interval $I$ and $f'(x)\ge0$ for all interior points of $I$, then $f$ is nondecreasing on $I$. If $f'(x)>0$, then $f$ is strictly increasing.

\end{corollary}

\meta{Proof ID}{\pidraw{COR-MONOTONICITY}}
\meta{Object type}{COROLLARY}
\meta{Immediate dependencies}{\pidref{THM-MVT}}
\meta{Additional transitive dependencies}{\pidref{AX-EVT}, \pidref{AX-REAL-FIELD}, \pidref{DEF-DERIVATIVE}, \pidref{THM-FERMAT}, \pidref{THM-ROLLE}}
\meta{Source traceability}{\textbf{Mathematical source:} source calculus framework, Sec. 3.3}

\begin{proofstrategy}Direct proof using the Mean Value Theorem.\end{proofstrategy}

\begin{proof}

Take $x_1<x_2$ in $I$. By the Mean Value Theorem there is $c\in(x_1,x_2)$ such that
\[
f(x_2)-f(x_1)=f'(c)(x_2-x_1).
\]
Since $x_2-x_1>0$, the sign of the difference is the sign of $f'(c)$.
Thus $f(x_2)\ge f(x_1)$ when $f'\ge0$, and the inequality is strict when
$f'>0$.

\end{proof}

\begin{proofanalysis}The MVT turns a local derivative sign into a finite-interval function comparison.\end{proofanalysis}

\begin{computationalrelevance}Monotonicity brackets roots and minima and justifies one-dimensional search logic.\end{computationalrelevance}

\subsubsection{Second derivative test}

\begin{theorem}[Second derivative test]\label{prf:THM-SECOND-DERIVATIVE-TEST}

Suppose $f''$ is continuous near $c$, $f'(c)=0$, and $f''(c)>0$. Then $f$ has a strict local minimum at $c$. If $f''(c)<0$, then $f$ has a strict local maximum.

\end{theorem}

\meta{Proof ID}{\pidraw{THM-SECOND-DERIVATIVE-TEST}}
\meta{Object type}{THEOREM}
\meta{Immediate dependencies}{\pidref{THM-MVT}, \pidref{COR-MONOTONICITY}}
\meta{Additional transitive dependencies}{\pidref{AX-EVT}, \pidref{AX-REAL-FIELD}, \pidref{DEF-DERIVATIVE}, \pidref{THM-FERMAT}, \pidref{THM-ROLLE}}
\meta{Source traceability}{\textbf{Mathematical source:} source calculus framework, Sec. 3.4}

\begin{proofstrategy}Direct proof using continuity and two applications of derivative-sign reasoning.\end{proofstrategy}

\begin{proof}

Assume $f''(c)>0$. Continuity of $f''$ gives $\eta>0$ such that
$f''(x)>0$ on $(c-\eta,c+\eta)$. Hence $f'$ is strictly increasing there by
\pid{COR-MONOTONICITY}. Because $f'(c)=0$, we have
$f'(x)<0$ for $x<c$ near $c$ and $f'(x)>0$ for $x>c$ near $c$.
Applying the monotonicity conclusion to $f$, the function decreases up to $c$
and increases after $c$. Therefore $f(c)$ is a strict local minimum.
The case $f''(c)<0$ follows by applying the argument to $-f$.

\end{proof}

\begin{proofanalysis}The sign of the second derivative controls monotonicity of the first derivative, which controls monotonicity of the function.\end{proofanalysis}

\begin{computationalrelevance}Curvature tests classify stationary points in engineering and economic optimization models.\end{computationalrelevance}

\subsection{Integral Calculus}

\subsubsection{Fundamental Theorem of Calculus, accumulation form}

\begin{theorem}[Fundamental Theorem of Calculus, accumulation form]\label{prf:THM-FTC-I}

If $f$ is continuous on $[a,b]$ and
\[
F(x)=\int_a^x f(t)\,dt,
\]
then $F$ is differentiable on $(a,b)$ and $F'(x)=f(x)$.

\end{theorem}

\meta{Proof ID}{\pidraw{THM-FTC-I}}
\meta{Object type}{THEOREM}
\meta{Immediate dependencies}{\pidref{DEF-RIEMANN-INTEGRAL}, \pidref{DEF-CONTINUITY}, \pidref{AX-RIEMANN-BASIC}}
\meta{Additional transitive dependencies}{\emph{None}}
\meta{Source traceability}{\textbf{Proof method:} Mathematical Proofs, direct epsilon-delta methodology; \textbf{Mathematical source:} source calculus framework, Sec. 4.4}

\begin{proofstrategy}Direct epsilon estimate on the difference quotient.\end{proofstrategy}

\begin{proof}

Fix $x\in(a,b)$. For sufficiently small nonzero $h$,
\[
\frac{F(x+h)-F(x)}{h}
=
\frac1h\int_x^{x+h}f(t)\,dt.
\]
Subtract $f(x)$:
\[
\frac{F(x+h)-F(x)}{h}-f(x)
=
\frac1h\int_x^{x+h}(f(t)-f(x))\,dt.
\]
Given $\varepsilon>0$, continuity of $f$ at $x$ gives $\eta>0$ such that
$|t-x|<\eta$ implies $|f(t)-f(x)|<\varepsilon$. For $0<|h|<\eta$, the
integral order estimate yields
\[
\left|
\frac1h\int_x^{x+h}(f(t)-f(x))\,dt
\right|
\le \varepsilon.
\]
Hence the difference quotient tends to $f(x)$.

\end{proof}

\begin{proofanalysis}The derivative of accumulated area is the average value of $f$ over a shrinking interval, which continuity forces toward the point value.\end{proofanalysis}

\begin{computationalrelevance}This theorem connects analytical antiderivatives to numerical quadrature and accumulation models.\end{computationalrelevance}

\subsubsection{Fundamental Theorem of Calculus, evaluation form}

\begin{theorem}[Fundamental Theorem of Calculus, evaluation form]\label{prf:THM-FTC-II}

If $F'$ is continuous on $[a,b]$, then
\[
\int_a^b F'(x)\,dx=F(b)-F(a).
\]

\end{theorem}

\meta{Proof ID}{\pidraw{THM-FTC-II}}
\meta{Object type}{THEOREM}
\meta{Immediate dependencies}{\pidref{THM-FTC-I}, \pidref{THM-MVT}}
\meta{Additional transitive dependencies}{\pidref{AX-EVT}, \pidref{AX-REAL-FIELD}, \pidref{AX-RIEMANN-BASIC}, \pidref{DEF-CONTINUITY}, \pidref{DEF-DERIVATIVE}, \pidref{DEF-RIEMANN-INTEGRAL}, \pidref{THM-FERMAT}, \pidref{THM-ROLLE}}
\meta{Source traceability}{\textbf{Mathematical source:} source calculus framework, Sec. 4.4}

\begin{proofstrategy}Direct proof by comparing two functions with the same derivative.\end{proofstrategy}

\begin{proof}

Let
\[
G(x)=\int_a^x F'(t)\,dt.
\]
By \pid{THM-FTC-I}, $G'(x)=F'(x)$ on $(a,b)$. Thus the function
$H(x)=F(x)-G(x)$ satisfies $H'(x)=0$ on $(a,b)$. By the Mean Value Theorem,
a differentiable function with zero derivative on an interval is constant, so
$H(b)=H(a)$. Since $G(a)=0$,
\[
F(b)-G(b)=F(a),
\]
and therefore
\[
G(b)=F(b)-F(a).
\]
But $G(b)=\int_a^b F'(x)\,dx$.

\end{proof}

\begin{proofanalysis}The proof reduces integral evaluation to uniqueness, up to an additive constant, of antiderivatives.\end{proofanalysis}

\begin{computationalrelevance}Closed-form integration and many verification strategies for numerical quadrature depend on this identity.\end{computationalrelevance}

\subsubsection{Substitution formula for definite integrals}

\begin{proposition}[Substitution formula for definite integrals]\label{prf:DER-SUBSTITUTION}

Let $g$ be continuously differentiable on $[\alpha,\beta]$ and let $f$ be continuous on an interval containing $g([\alpha,\beta])$. Then
\[
\int_\alpha^\beta f(g(t))g'(t)\,dt
=
\int_{g(\alpha)}^{g(\beta)} f(u)\,du.
\]

\end{proposition}

\meta{Proof ID}{\pidraw{DER-SUBSTITUTION}}
\meta{Object type}{DERIVATION}
\meta{Immediate dependencies}{\pidref{THM-CHAIN-RULE}, \pidref{THM-FTC-I}, \pidref{THM-FTC-II}}
\meta{Additional transitive dependencies}{\pidref{AX-EVT}, \pidref{AX-REAL-FIELD}, \pidref{AX-RIEMANN-BASIC}, \pidref{DEF-CONTINUITY}, \pidref{DEF-DERIVATIVE}, \pidref{DEF-LIMIT-FUNCTION}, \pidref{DEF-RIEMANN-INTEGRAL}, \pidref{THM-FERMAT}, \pidref{LEM-LIMIT-BOUNDED}, \pidref{THM-LIMIT-SUM}, \pidref{THM-ROLLE}, \pidref{THM-LIMIT-PRODUCT}, \pidref{THM-MVT}, \pidref{THM-DIFFERENTIABLE-CONTINUOUS}}
\meta{Source traceability}{\textbf{Mathematical source:} source calculus framework, Sec. 4.5; Sec. 8; \textbf{Numerical source:} Numerical Recipes, Sec. 4.5}
\meta{Catalog method IDs}{M16,M17}

\begin{proofstrategy}Algebraic derivation from the chain rule and the Fundamental Theorem.\end{proofstrategy}

\begin{proof}

Choose an antiderivative $F$ of $f$ on the relevant interval. By the chain rule,
\[
\frac{d}{dt}F(g(t))=f(g(t))g'(t).
\]
Applying \pid{THM-FTC-II},
\[
\int_\alpha^\beta f(g(t))g'(t)\,dt
=F(g(\beta))-F(g(\alpha)).
\]
A second application of the Fundamental Theorem gives
\[
F(g(\beta))-F(g(\alpha))
=\int_{g(\alpha)}^{g(\beta)} f(u)\,du.
\]

\end{proof}

\begin{proofanalysis}Substitution is the integral form of the chain rule.\end{proofanalysis}

\begin{computationalrelevance}Variable transformations are central to improper quadrature, coordinate changes, and numerical conditioning.\end{computationalrelevance}

\subsubsection{Integration by parts}

\begin{proposition}[Integration by parts]\label{prf:DER-INTEGRATION-BY-PARTS}

If $u$ and $v$ have continuous derivatives on $[a,b]$, then
\[
\int_a^b u(x)v'(x)\,dx
=
[u(x)v(x)]_a^b-\int_a^b u'(x)v(x)\,dx.
\]

\end{proposition}

\meta{Proof ID}{\pidraw{DER-INTEGRATION-BY-PARTS}}
\meta{Object type}{DERIVATION}
\meta{Immediate dependencies}{\pidref{THM-PRODUCT-RULE}, \pidref{THM-FTC-II}}
\meta{Additional transitive dependencies}{\pidref{AX-EVT}, \pidref{AX-REAL-FIELD}, \pidref{AX-RIEMANN-BASIC}, \pidref{DEF-CONTINUITY}, \pidref{DEF-DERIVATIVE}, \pidref{DEF-LIMIT-FUNCTION}, \pidref{DEF-RIEMANN-INTEGRAL}, \pidref{THM-FERMAT}, \pidref{LEM-LIMIT-BOUNDED}, \pidref{THM-LIMIT-SUM}, \pidref{THM-FTC-I}, \pidref{THM-ROLLE}, \pidref{THM-LIMIT-PRODUCT}, \pidref{THM-MVT}, \pidref{THM-DIFFERENTIABLE-CONTINUOUS}}
\meta{Source traceability}{\textbf{Mathematical source:} source calculus framework, Sec. 8.2}

\begin{proofstrategy}Algebraic derivation from the product rule.\end{proofstrategy}

\begin{proof}

The product rule gives
\[
(uv)'=u'v+uv'.
\]
Integrating both sides over $[a,b]$ and applying \pid{THM-FTC-II},
\[
u(b)v(b)-u(a)v(a)
=
\int_a^b u'v\,dx+\int_a^b uv'\,dx.
\]
Rearranging yields the formula.

\end{proof}

\begin{proofanalysis}The identity transfers a derivative from one factor to the other.\end{proofanalysis}

\begin{computationalrelevance}Integration by parts underlies analytical benchmark integrals and many asymptotic and Fourier calculations.\end{computationalrelevance}

\subsection{Sequences and Series}

\subsubsection{Geometric series formula}

\begin{theorem}[Geometric series formula]\label{prf:THM-GEOMETRIC-SERIES}

For $r\in\mathbb{R}$ with $|r|<1$,
\[
\sum_{k=0}^{\infty} r^k=\frac{1}{1-r}.
\]

\end{theorem}

\meta{Proof ID}{\pidraw{THM-GEOMETRIC-SERIES}}
\meta{Object type}{THEOREM}
\meta{Immediate dependencies}{\pidref{DEF-SERIES}, \pidref{AX-REAL-FIELD}, \pidref{AX-ELEMENTARY-SEQUENCE-LIMITS}}
\meta{Additional transitive dependencies}{\emph{None}}
\meta{Source traceability}{\textbf{Proof method:} Mathematical Proofs, Secs. 14.1--14.2; \textbf{Mathematical source:} source calculus framework, Secs. 9.1--9.2}

\begin{proofstrategy}Direct proof through partial sums and a limit.\end{proofstrategy}

\begin{proof}

For $r\ne1$, the finite partial sum satisfies
\[
s_n=1+r+\cdots+r^n=\frac{1-r^{n+1}}{1-r},
\]
because $(1-r)s_n=1-r^{n+1}$. If $|r|<1$, then $r^{n+1}\to0$.
Therefore
\[
s_n\to\frac{1}{1-r},
\]
which is exactly the definition of convergence of the series.

\end{proof}

\begin{proofanalysis}The infinite statement is reduced to a finite algebraic identity plus one sequence limit.\end{proofanalysis}

\begin{computationalrelevance}Geometric convergence models iterative error reduction, discounting, and series-based approximations.\end{computationalrelevance}

\subsubsection{Convergence of a series forces its terms to zero}

\begin{theorem}[Convergence of a series forces its terms to zero]\label{prf:THM-SERIES-TERM-TEST}

If $\sum_{k=1}^{\infty}a_k$ converges, then $a_n\to0$.

\end{theorem}

\meta{Proof ID}{\pidraw{THM-SERIES-TERM-TEST}}
\meta{Object type}{THEOREM}
\meta{Immediate dependencies}{\pidref{DEF-SERIES}, \pidref{THM-LIMIT-SUM}}
\meta{Additional transitive dependencies}{\pidref{DEF-LIMIT-FUNCTION}}
\meta{Source traceability}{\textbf{Proof method:} Mathematical Proofs, Exercise 14.16 context; \textbf{Mathematical source:} source calculus framework, Sec. 9.2}

\begin{proofstrategy}Direct proof from consecutive partial sums.\end{proofstrategy}

\begin{proof}

Let $s_n=\sum_{k=1}^{n}a_k$ and suppose $s_n\to s$. For $n\ge2$,
\[
a_n=s_n-s_{n-1}.
\]
Both $s_n$ and the shifted sequence $s_{n-1}$ tend to $s$. Hence
\[
a_n\to s-s=0
\]
by the limit laws.

\end{proof}

\begin{proofanalysis}Each term is the difference between two consecutive approximations to the same limiting sum.\end{proofanalysis}

\begin{computationalrelevance}This is a necessary diagnostic for any series used in function approximation or convergence acceleration.\end{computationalrelevance}

\subsubsection{Taylor theorem with Lagrange remainder}

\begin{theorem}[Taylor theorem with Lagrange remainder]\label{prf:THM-TAYLOR-LAGRANGE}

Suppose $f$ has $n+1$ derivatives on an interval containing $a$ and $x$, with the needed continuity to apply Rolle's theorem. Then for some $\xi$ between $a$ and $x$,
\[
f(x)=\sum_{k=0}^{n}\frac{f^{(k)}(a)}{k!}(x-a)^k
+\frac{f^{(n+1)}(\xi)}{(n+1)!}(x-a)^{n+1}.
\]

\end{theorem}

\meta{Proof ID}{\pidraw{THM-TAYLOR-LAGRANGE}}
\meta{Object type}{THEOREM}
\meta{Immediate dependencies}{\pidref{THM-ROLLE}}
\meta{Additional transitive dependencies}{\pidref{AX-EVT}, \pidref{AX-REAL-FIELD}, \pidref{DEF-DERIVATIVE}, \pidref{THM-FERMAT}}
\meta{Source traceability}{\textbf{Proof method:} Mathematical Proofs, existence proof plus repeated Rolle-style reasoning; \textbf{Mathematical source:} source calculus framework, Secs. 9.7--9.10}

\begin{proofstrategy}Existence proof by constructing an auxiliary function and applying Rolle's theorem repeatedly.\end{proofstrategy}

\begin{proof}

Let
\[
P_n(t)=\sum_{k=0}^{n}\frac{f^{(k)}(a)}{k!}(t-a)^k,
\qquad
R=f(x)-P_n(x).
\]
If $x=a$ the result is immediate, so assume $x\ne a$ and set
\[
K=\frac{R}{(x-a)^{n+1}}.
\]
Define
\[
\phi(t)=f(t)-P_n(t)-K(t-a)^{n+1}.
\]
Then $\phi(x)=0$, while
\[
\phi(a)=\phi'(a)=\cdots=\phi^{(n)}(a)=0.
\]
Apply Rolle's theorem to $\phi$ on the interval with endpoints $a$ and $x$,
then to $\phi'$, and continue. After $n+1$ applications there is a point
$\xi$ between $a$ and $x$ with $\phi^{(n+1)}(\xi)=0$. Since
\[
\phi^{(n+1)}(t)=f^{(n+1)}(t)-(n+1)!K,
\]
we have
\[
K=\frac{f^{(n+1)}(\xi)}{(n+1)!}.
\]
Substituting the definition of $K$ yields the stated formula.

\end{proof}

\begin{proofanalysis}The auxiliary polynomial is chosen so that many derivatives vanish at one endpoint and the residual vanishes at the other; repeated Rolle steps force the remainder coefficient.\end{proofanalysis}

\begin{computationalrelevance}Taylor remainders justify finite-difference errors, Newton's local rate, Euler truncation error, and many quadrature formulas.\end{computationalrelevance}

\subsection{Parametric and Vector Calculus}

\subsubsection{Derivative of a parametrically defined curve}

\begin{proposition}[Derivative of a parametrically defined curve]\label{prf:DER-PARAMETRIC-DERIVATIVE}

If $x=x(t)$ and $y=y(t)$ are differentiable at $t_0$ with $x'(t_0)\ne0$, and locally $y$ can be regarded as a differentiable function of $x$, then
\[
\frac{dy}{dx}\Big|_{t_0}=\frac{y'(t_0)}{x'(t_0)}.
\]

\end{proposition}

\meta{Proof ID}{\pidraw{DER-PARAMETRIC-DERIVATIVE}}
\meta{Object type}{DERIVATION}
\meta{Immediate dependencies}{\pidref{THM-CHAIN-RULE}}
\meta{Additional transitive dependencies}{\pidref{AX-REAL-FIELD}, \pidref{DEF-CONTINUITY}, \pidref{DEF-DERIVATIVE}, \pidref{DEF-LIMIT-FUNCTION}, \pidref{LEM-LIMIT-BOUNDED}, \pidref{THM-LIMIT-SUM}, \pidref{THM-LIMIT-PRODUCT}, \pidref{THM-DIFFERENTIABLE-CONTINUOUS}}
\meta{Source traceability}{\textbf{Mathematical source:} source calculus framework, Sec. 10.3}

\begin{proofstrategy}Algebraic derivation from the chain rule.\end{proofstrategy}

\begin{proof}

Write $y(t)=Y(x(t))$ near $t_0$. The chain rule gives
\[
y'(t_0)=Y'(x(t_0))\,x'(t_0).
\]
Since $x'(t_0)\ne0$,
\[
Y'(x(t_0))=\frac{y'(t_0)}{x'(t_0)}.
\]
This derivative is $dy/dx$ along the curve.

\end{proof}

\begin{proofanalysis}The formula is valid only where the parameter actually changes the $x$ coordinate to first order.\end{proofanalysis}

\begin{computationalrelevance}Parametric derivatives drive velocity/acceleration models, curve geometry, and trajectory optimization.\end{computationalrelevance}

\subsection{Multivariable Calculus}

\subsubsection{Chain rule along a differentiable curve}

\begin{theorem}[Chain rule along a differentiable curve]\label{prf:THM-CURVE-CHAIN}

If $F:\mathbb{R}^n\to\mathbb{R}$ is differentiable at $r(t_0)$ and
$r:\mathbb{R}\to\mathbb{R}^n$ is differentiable at $t_0$, then
\[
\frac{d}{dt}F(r(t))\Big|_{t=t_0}
=
\nabla F(r(t_0))\cdot r'(t_0).
\]

\end{theorem}

\meta{Proof ID}{\pidraw{THM-CURVE-CHAIN}}
\meta{Object type}{THEOREM}
\meta{Immediate dependencies}{\pidref{DEF-DIFFERENTIABLE-RN}}
\meta{Additional transitive dependencies}{\emph{None}}
\meta{Source traceability}{\textbf{Mathematical source:} source calculus framework, Secs. 13.5 and 15.3}

\begin{proofstrategy}Direct proof by composing the linear approximation with the curve increment.\end{proofstrategy}

\begin{proof}

Let $a=r(t_0)$ and $\Delta r=r(t_0+h)-r(t_0)$. Differentiability of $F$ gives
\[
F(a+\Delta r)-F(a)=\nabla F(a)\cdot\Delta r+o(\|\Delta r\|).
\]
Divide by $h$. Since $r$ is differentiable,
$\Delta r/h\to r'(t_0)$ and $\|\Delta r\|/|h|$ remains bounded near $0$.
Therefore
\[
\frac{o(\|\Delta r\|)}{h}\to0,
\]
and the limit is
\[
\nabla F(a)\cdot r'(t_0).
\]

\end{proof}

\begin{proofanalysis}The multivariable linear approximation is evaluated on the one-dimensional motion of the curve.\end{proofanalysis}

\begin{computationalrelevance}This identity is the engine behind line-integral potential theorems and field-based dynamics.\end{computationalrelevance}

\subsubsection{Directional derivative equals gradient dot direction}

\begin{theorem}[Directional derivative equals gradient dot direction]\label{prf:THM-DIRECTIONAL-GRADIENT}

If $F:\mathbb{R}^n\to\mathbb{R}$ is differentiable at $a$, then for every unit vector $u$,
\[
D_uF(a)=\nabla F(a)\cdot u.
\]

\end{theorem}

\meta{Proof ID}{\pidraw{THM-DIRECTIONAL-GRADIENT}}
\meta{Object type}{THEOREM}
\meta{Immediate dependencies}{\pidref{DEF-DIFFERENTIABLE-RN}, \pidref{DEF-DIRECTIONAL-DERIVATIVE}}
\meta{Additional transitive dependencies}{\emph{None}}
\meta{Source traceability}{\textbf{Mathematical source:} source calculus framework, Sec. 13.6}

\begin{proofstrategy}Direct proof from the definition of differentiability.\end{proofstrategy}

\begin{proof}

Differentiability gives
\[
F(a+h)=F(a)+\nabla F(a)\cdot h+r(h),
\qquad
\frac{r(h)}{\|h\|}\to0.
\]
Set $h=tu$. For $t\ne0$,
\[
\frac{F(a+tu)-F(a)}{t}
=
\nabla F(a)\cdot u+\frac{r(tu)}{t}.
\]
Because $\|u\|=1$,
\[
\left|\frac{r(tu)}{t}\right|
=
\frac{|r(tu)|}{\|tu\|}\to0.
\]
Taking the limit proves the formula.

\end{proof}

\begin{proofanalysis}Differentiability is stronger than existence of partial derivatives: it gives one linear map that controls every direction.\end{proofanalysis}

\begin{computationalrelevance}The gradient is the computational object used in multivariable optimization, sensitivity analysis, and gradient-driven trajectories.\end{computationalrelevance}

\subsection{Vector Analysis}

\subsubsection{Fundamental Theorem of Line Integrals}

\begin{theorem}[Fundamental Theorem of Line Integrals]\label{prf:THM-FUNDAMENTAL-LINE-INTEGRAL}

Let $F=\nabla\phi$ on a region containing a piecewise smooth curve $C$ from $A$ to $B$. Then
\[
\int_C F\cdot dr=\phi(B)-\phi(A).
\]

\end{theorem}

\meta{Proof ID}{\pidraw{THM-FUNDAMENTAL-LINE-INTEGRAL}}
\meta{Object type}{THEOREM}
\meta{Immediate dependencies}{\pidref{THM-CURVE-CHAIN}, \pidref{THM-FTC-II}}
\meta{Additional transitive dependencies}{\pidref{AX-EVT}, \pidref{AX-REAL-FIELD}, \pidref{AX-RIEMANN-BASIC}, \pidref{DEF-CONTINUITY}, \pidref{DEF-DERIVATIVE}, \pidref{DEF-DIFFERENTIABLE-RN}, \pidref{DEF-RIEMANN-INTEGRAL}, \pidref{THM-FERMAT}, \pidref{THM-FTC-I}, \pidref{THM-ROLLE}, \pidref{THM-MVT}}
\meta{Source traceability}{\textbf{Mathematical source:} source calculus framework, Sec. 15.3}

\begin{proofstrategy}Direct reduction of the line integral to a one-variable derivative.\end{proofstrategy}

\begin{proof}

Let $r:[\alpha,\beta]\to\mathbb{R}^n$ parametrize $C$ with
$r(\alpha)=A$ and $r(\beta)=B$. Then
\[
\int_C F\cdot dr
=
\int_\alpha^\beta \nabla\phi(r(t))\cdot r'(t)\,dt.
\]
By \pid{THM-CURVE-CHAIN}, the integrand equals
$\frac{d}{dt}\phi(r(t))$. Therefore \pid{THM-FTC-II} gives
\[
\int_C F\cdot dr
=
\phi(r(\beta))-\phi(r(\alpha))
=
\phi(B)-\phi(A).
\]

\end{proof}

\begin{proofanalysis}A conservative vector field converts a path integral into an endpoint calculation.\end{proofanalysis}

\begin{computationalrelevance}Potential fields, work, gravitational/electric models, and path-independent energy computations use this theorem.\end{computationalrelevance}

\subsubsection{Conservative fields are path independent}

\begin{corollary}[Conservative fields are path independent]\label{prf:COR-PATH-INDEPENDENCE}

If $F=\nabla\phi$ on a region and $C_1,C_2$ are two piecewise smooth curves with the same start and end points, then
\[
\int_{C_1}F\cdot dr=\int_{C_2}F\cdot dr.
\]

\end{corollary}

\meta{Proof ID}{\pidraw{COR-PATH-INDEPENDENCE}}
\meta{Object type}{COROLLARY}
\meta{Immediate dependencies}{\pidref{THM-FUNDAMENTAL-LINE-INTEGRAL}}
\meta{Additional transitive dependencies}{\pidref{AX-EVT}, \pidref{AX-REAL-FIELD}, \pidref{AX-RIEMANN-BASIC}, \pidref{DEF-CONTINUITY}, \pidref{DEF-DERIVATIVE}, \pidref{DEF-DIFFERENTIABLE-RN}, \pidref{DEF-RIEMANN-INTEGRAL}, \pidref{THM-FERMAT}, \pidref{THM-CURVE-CHAIN}, \pidref{THM-FTC-I}, \pidref{THM-ROLLE}, \pidref{THM-MVT}, \pidref{THM-FTC-II}}
\meta{Source traceability}{\textbf{Mathematical source:} source calculus framework, Sec. 15.3}

\begin{proofstrategy}Direct proof from endpoint evaluation.\end{proofstrategy}

\begin{proof}

Apply \pid{THM-FUNDAMENTAL-LINE-INTEGRAL} to each curve. Both integrals equal
$\phi(B)-\phi(A)$, where $A$ and $B$ are the common endpoints. Hence the two
line integrals are equal.

\end{proof}

\begin{proofanalysis}The path disappears from the formula once a potential exists.\end{proofanalysis}

\begin{computationalrelevance}This is the mathematical basis for replacing expensive path integration by endpoint evaluation in conservative fields.\end{computationalrelevance}

\subsection{Numerical Approximation}

\subsubsection{A nonzero polynomial has at most its degree many distinct real roots}

\begin{lemma}[A nonzero polynomial has at most its degree many distinct real roots]\label{prf:LEM-POLYNOMIAL-ROOT-COUNT}

A nonzero polynomial of degree $n$ has at most $n$ distinct real roots.

\end{lemma}

\meta{Proof ID}{\pidraw{LEM-POLYNOMIAL-ROOT-COUNT}}
\meta{Object type}{LEMMA}
\meta{Immediate dependencies}{\pidref{AX-POLY-FACTOR}}
\meta{Additional transitive dependencies}{\emph{None}}
\meta{Source traceability}{\textbf{Proof method:} Mathematical Proofs, Ch. 6 induction methodology}

\begin{proofstrategy}Mathematical induction on the degree.\end{proofstrategy}

\begin{proof}

For degree $0$, a nonzero constant has no roots. Assume the statement for degree
$n-1$, and let $p$ be a nonzero degree-$n$ polynomial. If $p$ has no root there is
nothing to prove. If $r$ is a root, the factor theorem gives
$p(x)=(x-r)q(x)$ with $\deg q=n-1$. Every root of $p$ distinct from $r$ is a
root of $q$, which by the induction hypothesis has at most $n-1$ distinct roots.
Thus $p$ has at most $n$ distinct roots.

\end{proof}

\begin{proofanalysis}Each distinct root consumes at least one linear factor.\end{proofanalysis}

\begin{computationalrelevance}Root counting proves uniqueness of interpolation polynomials and supports polynomial root reasoning.\end{computationalrelevance}

\subsubsection{Uniqueness of polynomial interpolation}

\begin{theorem}[Uniqueness of polynomial interpolation]\label{prf:THM-INTERPOLATION-UNIQUENESS}

Given $n+1$ distinct nodes $x_0,\ldots,x_n$ and values $y_0,\ldots,y_n$, there is at most one polynomial of degree at most $n$ satisfying $p(x_i)=y_i$ for all $i$.

\end{theorem}

\meta{Proof ID}{\pidraw{THM-INTERPOLATION-UNIQUENESS}}
\meta{Object type}{THEOREM}
\meta{Immediate dependencies}{\pidref{LEM-POLYNOMIAL-ROOT-COUNT}}
\meta{Additional transitive dependencies}{\pidref{AX-POLY-FACTOR}}
\meta{Source traceability}{\textbf{Numerical source:} Numerical Recipes, Sec. 3.2}
\meta{Catalog method IDs}{M02}

\begin{proofstrategy}Proof by contradiction through a difference polynomial.\end{proofstrategy}

\begin{proof}

Suppose $p$ and $q$ are two such polynomials. Their difference
$r=p-q$ has degree at most $n$ and satisfies
$r(x_i)=0$ for all $n+1$ distinct nodes. By
\pid{LEM-POLYNOMIAL-ROOT-COUNT}, a nonzero degree-at-most-$n$ polynomial cannot
have $n+1$ distinct roots. Hence $r$ is the zero polynomial and $p=q$.

\end{proof}

\begin{proofanalysis}Interpolation uniqueness is a root-counting statement in disguise.\end{proofanalysis}

\begin{computationalrelevance}It guarantees that polynomial interpolation routines are computing a mathematically well-defined object.\end{computationalrelevance}

\subsubsection{Lagrange interpolation formula}

\begin{theorem}[Lagrange interpolation formula]\label{prf:THM-LAGRANGE-INTERPOLATION}

For distinct nodes $x_0,\ldots,x_n$, the polynomial
\[
p(x)=\sum_{i=0}^n y_i
\prod_{\substack{0\le j\le n\\j\ne i}}
\frac{x-x_j}{x_i-x_j}
\]
is the unique polynomial of degree at most $n$ satisfying $p(x_i)=y_i$.

\end{theorem}

\meta{Proof ID}{\pidraw{THM-LAGRANGE-INTERPOLATION}}
\meta{Object type}{THEOREM}
\meta{Immediate dependencies}{\pidref{THM-INTERPOLATION-UNIQUENESS}}
\meta{Additional transitive dependencies}{\pidref{AX-POLY-FACTOR}, \pidref{LEM-POLYNOMIAL-ROOT-COUNT}}
\meta{Source traceability}{\textbf{Numerical source:} Numerical Recipes, Secs. 3.2 and 3.5}
\meta{Catalog method IDs}{M02}

\begin{proofstrategy}Constructive existence proof followed by uniqueness.\end{proofstrategy}

\begin{proof}

Define the basis polynomials
\[
L_i(x)=
\prod_{j\ne i}\frac{x-x_j}{x_i-x_j}.
\]
Distinctness of the nodes makes every denominator nonzero. At a node $x_k$,
$L_i(x_k)=0$ for $i\ne k$ and $L_k(x_k)=1$. Therefore
\[
p(x_k)=\sum_{i=0}^n y_i L_i(x_k)=y_k.
\]
The degree is at most $n$. Uniqueness follows from
\pid{THM-INTERPOLATION-UNIQUENESS}.

\end{proof}

\begin{proofanalysis}The basis functions are cardinal: each isolates exactly one data value at the nodes.\end{proofanalysis}

\begin{computationalrelevance}This is the exact polynomial object approximated or reorganized by practical interpolation algorithms.\end{computationalrelevance}

\subsubsection{Second-order error of the centered difference}

\begin{theorem}[Second-order error of the centered difference]\label{prf:ERR-CENTERED-DIFFERENCE}

If $f$ has a continuous third derivative on $[x-h,x+h]$, then
\[
\frac{f(x+h)-f(x-h)}{2h}=f'(x)+E_h,
\qquad
|E_h|\le \frac{h^2}{6}\max_{|t-x|\le|h|}|f'''(t)|.
\]
Hence the centered difference is second-order accurate.

\end{theorem}

\meta{Proof ID}{\pidraw{ERR-CENTERED-DIFFERENCE}}
\meta{Object type}{ERROR-THEOREM}
\meta{Immediate dependencies}{\pidref{THM-TAYLOR-LAGRANGE}, \pidref{DEF-BIG-O}}
\meta{Additional transitive dependencies}{\pidref{AX-EVT}, \pidref{AX-REAL-FIELD}, \pidref{DEF-DERIVATIVE}, \pidref{THM-FERMAT}, \pidref{THM-ROLLE}}
\meta{Source traceability}{\textbf{Numerical source:} Numerical Recipes, Sec. 5.7}
\meta{Catalog method IDs}{M01}

\begin{proofstrategy}Direct proof by symmetric Taylor expansions.\end{proofstrategy}

\begin{proof}

Taylor expansion about $x$ gives
\[
f(x+h)=f(x)+hf'(x)+\frac{h^2}{2}f''(x)+\frac{h^3}{6}f'''(\xi_+),
\]
\[
f(x-h)=f(x)-hf'(x)+\frac{h^2}{2}f''(x)-\frac{h^3}{6}f'''(\xi_-),
\]
for points $\xi_\pm$ between $x$ and $x\pm h$. Subtracting and dividing by
$2h$ yields
\[
\frac{f(x+h)-f(x-h)}{2h}
=
f'(x)+\frac{h^2}{12}\bigl(f'''(\xi_+)+f'''(\xi_-)\bigr).
\]
The displayed bound follows immediately, and therefore the error is $O(h^2)$.

\end{proof}

\begin{proofanalysis}Symmetry cancels the even Taylor terms, raising the accuracy order relative to a one-sided difference.\end{proofanalysis}

\begin{computationalrelevance}This explains the higher truncation accuracy of centered numerical differentiation used in scientific computing.\end{computationalrelevance}

\subsubsection{First-order error of the forward difference}

\begin{theorem}[First-order error of the forward difference]\label{prf:ERR-FORWARD-DIFFERENCE}

If $f$ has a continuous second derivative between $x$ and $x+h$, then for some $\xi$ between them,
\[
\frac{f(x+h)-f(x)}{h}
=
f'(x)+\frac{h}{2}f''(\xi).
\]
Consequently the forward-difference truncation error is $O(h)$.

\end{theorem}

\meta{Proof ID}{\pidraw{ERR-FORWARD-DIFFERENCE}}
\meta{Object type}{ERROR-THEOREM}
\meta{Immediate dependencies}{\pidref{THM-TAYLOR-LAGRANGE}, \pidref{DEF-BIG-O}}
\meta{Additional transitive dependencies}{\pidref{AX-EVT}, \pidref{AX-REAL-FIELD}, \pidref{DEF-DERIVATIVE}, \pidref{THM-FERMAT}, \pidref{THM-ROLLE}}
\meta{Source traceability}{\textbf{Numerical source:} Numerical Recipes, Sec. 5.7}
\meta{Catalog method IDs}{M01}

\begin{proofstrategy}Direct Taylor expansion with remainder.\end{proofstrategy}

\begin{proof}

Taylor's theorem of order $1$ gives
\[
f(x+h)=f(x)+hf'(x)+\frac{h^2}{2}f''(\xi)
\]
for some $\xi$ between $x$ and $x+h$. Subtract $f(x)$ and divide by $h\ne0$:
\[
\frac{f(x+h)-f(x)}{h}
=
f'(x)+\frac{h}{2}f''(\xi).
\]
If $|f''|\le M$ on a fixed neighborhood, the absolute truncation error is at most
$M|h|/2$, which is $O(h)$.

\end{proof}

\begin{proofanalysis}The first omitted Taylor term determines the leading truncation order.\end{proofanalysis}

\begin{computationalrelevance}This is the mathematical basis for one-sided numerical derivative routines and step-size/error reasoning.\end{computationalrelevance}

\subsection{Root Finding}

\subsubsection{Newton iteration from tangent-line linearization}

\begin{proposition}[Newton iteration from tangent-line linearization]\label{prf:DER-NEWTON-ITERATION}

For a differentiable $f$ with $f'(x_n)\ne0$, the zero of the tangent-line approximation at $x_n$ is
\[
x_{n+1}=x_n-\frac{f(x_n)}{f'(x_n)}.
\]

\end{proposition}

\meta{Proof ID}{\pidraw{DER-NEWTON-ITERATION}}
\meta{Object type}{DERIVATION}
\meta{Immediate dependencies}{\pidref{DEF-DERIVATIVE}}
\meta{Additional transitive dependencies}{\emph{None}}
\meta{Source traceability}{\textbf{Mathematical source:} source calculus framework, Sec. 3.8; \textbf{Numerical source:} Numerical Recipes, Sec. 9.4}
\meta{Catalog method IDs}{M07}

\begin{proofstrategy}Algebraic derivation from the tangent-line equation.\end{proofstrategy}

\begin{proof}

The tangent line to $y=f(x)$ at $x_n$ is
\[
y=f(x_n)+f'(x_n)(x-x_n).
\]
Set $y=0$ and solve for its $x$-intercept:
\[
0=f(x_n)+f'(x_n)(x_{n+1}-x_n),
\]
so, because $f'(x_n)\ne0$,
\[
x_{n+1}=x_n-\frac{f(x_n)}{f'(x_n)}.
\]

\end{proof}

\begin{proofanalysis}Newton's method is a local linearization method; the iteration formula is not itself a convergence theorem.\end{proofanalysis}

\begin{computationalrelevance}The distinction between derivation and convergence is essential when using Newton solvers in production code.\end{computationalrelevance}

\subsubsection{Opposite-sign bracket preserves a root}

\begin{proposition}[Opposite-sign bracket preserves a root]\label{prf:ALG-SIGN-BRACKET-INVARIANT}

Let $f$ be continuous on $[a,b]$ with $f(a)f(b)\le0$. If $c\in(a,b)$ is chosen and one replaces $[a,b]$ by a subinterval whose endpoints still satisfy a nonpositive sign product, then the new interval contains a zero of $f$.

\end{proposition}

\meta{Proof ID}{\pidraw{ALG-SIGN-BRACKET-INVARIANT}}
\meta{Object type}{ALGORITHM-INVARIANT}
\meta{Immediate dependencies}{\pidref{THM-IVT}}
\meta{Additional transitive dependencies}{\pidref{AX-REAL-COMPLETENESS}, \pidref{AX-REAL-FIELD}, \pidref{DEF-CONTINUITY}}
\meta{Source traceability}{\textbf{Numerical source:} Numerical Recipes, Secs. 9.1--9.3}
\meta{Catalog method IDs}{M04,M05,M06}

\begin{proofstrategy}Direct proof from the Intermediate Value Theorem.\end{proofstrategy}

\begin{proof}

Let the retained interval be $[\alpha,\beta]\subseteq[a,b]$ and assume
$f(\alpha)f(\beta)\le0$. If one endpoint value is zero, that endpoint is a root.
Otherwise the endpoint values have opposite signs, so $0$ lies between them.
By \pid{THM-IVT}, there is $r\in(\alpha,\beta)$ with $f(r)=0$.

\end{proof}

\begin{proofanalysis}The invariant is topological: the precise rule used to choose the interior trial point is irrelevant as long as the sign bracket is retained.\end{proofanalysis}

\begin{computationalrelevance}Bisection, false position, and safeguarded Brent methods all rely on this invariant.\end{computationalrelevance}

\subsubsection{Bisection preserves a sign bracket and halves its width}

\begin{proposition}[Bisection preserves a sign bracket and halves its width]\label{prf:ALG-BISECTION-INVARIANT}

Under the assumptions of a sign bracket, one bisection step selects a child interval that contains a root and has exactly half the parent width.

\end{proposition}

\meta{Proof ID}{\pidraw{ALG-BISECTION-INVARIANT}}
\meta{Object type}{ALGORITHM-INVARIANT}
\meta{Immediate dependencies}{\pidref{ALG-SIGN-BRACKET-INVARIANT}}
\meta{Additional transitive dependencies}{\pidref{AX-REAL-COMPLETENESS}, \pidref{AX-REAL-FIELD}, \pidref{DEF-CONTINUITY}, \pidref{THM-IVT}}
\meta{Source traceability}{\textbf{Numerical source:} Numerical Recipes, Sec. 9.1}
\meta{Catalog method IDs}{M04}

\begin{proofstrategy}Direct proof by cases.\end{proofstrategy}

\begin{proof}

Let $m=(a+b)/2$. If $f(m)=0$, the root has been found. Otherwise either
$f(a)f(m)<0$ or $f(m)f(b)<0$; retain the corresponding half interval.
By \pid{ALG-SIGN-BRACKET-INVARIANT}, the retained interval contains a root.
Its width is
\[
m-a=\frac{b-a}{2}
\quad\text{or}\quad
b-m=\frac{b-a}{2}.
\]

\end{proof}

\begin{proofanalysis}Bisection combines a qualitative existence invariant with an exact geometric contraction.\end{proofanalysis}

\begin{computationalrelevance}This is the correctness invariant of the standard bisection root solver.\end{computationalrelevance}

\subsubsection{Convergence of bisection}

\begin{theorem}[Convergence of bisection]\label{prf:CONV-BISECTION}

Let $f$ be continuous on $[a_0,b_0]$ with $f(a_0)f(b_0)\le0$. Repeated bisection produces nested intervals $I_n=[a_n,b_n]$ with
\[
|I_n|=\frac{b_0-a_0}{2^n}
\]
and a unique common point $r$ satisfying $f(r)=0$. Any sequence $x_n\in I_n$ converges to $r$.

\end{theorem}

\meta{Proof ID}{\pidraw{CONV-BISECTION}}
\meta{Object type}{CONVERGENCE-THEOREM}
\meta{Immediate dependencies}{\pidref{ALG-BISECTION-INVARIANT}, \pidref{THM-NESTED-INTERVAL}, \pidref{DEF-CONTINUITY}}
\meta{Additional transitive dependencies}{\pidref{AX-REAL-COMPLETENESS}, \pidref{AX-REAL-FIELD}, \pidref{THM-IVT}, \pidref{ALG-SIGN-BRACKET-INVARIANT}}
\meta{Source traceability}{\textbf{Numerical source:} Numerical Recipes, Sec. 9.1}
\meta{Catalog method IDs}{M04}

\begin{proofstrategy}Invariant plus limit argument.\end{proofstrategy}

\begin{proof}

By induction using \pid{ALG-BISECTION-INVARIANT}, each $I_n$ contains a root,
$I_{n+1}\subseteq I_n$, and
\[
b_n-a_n=\frac{b_0-a_0}{2^n}\to0.
\]
By \pid{THM-NESTED-INTERVAL}, the intervals have exactly one common point $r$.
For each $n$, choose a root $z_n\in I_n$, whose existence follows from the bracket
invariant. Since both $z_n$ and $r$ lie in $I_n$,
$|z_n-r|\le |I_n|\to0$, so $z_n\to r$. Continuity then gives
$f(r)=\lim f(z_n)=0$. Finally, for any $x_n\in I_n$,
$|x_n-r|\le|I_n|\to0$, hence $x_n\to r$.

\end{proof}

\begin{proofanalysis}The convergence proof separates existence of a root, geometric interval contraction, and continuity at the limiting point.\end{proofanalysis}

\begin{computationalrelevance}It provides a deterministic iteration bound: after $n$ steps the uncertainty width is $(b_0-a_0)/2^n$.\end{computationalrelevance}

\subsubsection{Local quadratic error estimate for Newton's method}

\begin{theorem}[Local quadratic error estimate for Newton's method]\label{prf:CONV-NEWTON-LOCAL}

Let $r$ be a simple root of $f$. Suppose $f$ has a continuous second derivative on an interval $I$ containing $r$, $|f'(x)|\ge m>0$ and $|f''(x)|\le M$ on $I$, and the Newton iterates remain in $I$. Then
\[
|x_{n+1}-r|\le \frac{M}{2m}|x_n-r|^2.
\]
Thus, once the iterates are in such a region, the error decreases quadratically.

\end{theorem}

\meta{Proof ID}{\pidraw{CONV-NEWTON-LOCAL}}
\meta{Object type}{CONVERGENCE-THEOREM}
\meta{Immediate dependencies}{\pidref{DER-NEWTON-ITERATION}, \pidref{THM-TAYLOR-LAGRANGE}}
\meta{Additional transitive dependencies}{\pidref{AX-EVT}, \pidref{AX-REAL-FIELD}, \pidref{DEF-DERIVATIVE}, \pidref{THM-FERMAT}, \pidref{THM-ROLLE}}
\meta{Source traceability}{\textbf{Numerical source:} Numerical Recipes, Sec. 9.4}
\meta{Catalog method IDs}{M07}

\begin{proofstrategy}Direct Taylor remainder estimate.\end{proofstrategy}

\begin{proof}

Apply Taylor's theorem to $f$ about $x_n$ and evaluate at the root $r$:
\[
0=f(r)=f(x_n)+f'(x_n)(r-x_n)
+\frac12 f''(\xi_n)(r-x_n)^2
\]
for some $\xi_n$ between $x_n$ and $r$. Rearranging,
\[
x_n-\frac{f(x_n)}{f'(x_n)}-r
=
\frac{f''(\xi_n)}{2f'(x_n)}(x_n-r)^2.
\]
By \pid{DER-NEWTON-ITERATION}, the left side is $x_{n+1}-r$. Therefore
\[
|x_{n+1}-r|
\le
\frac{M}{2m}|x_n-r|^2.
\]

\end{proof}

\begin{proofanalysis}Quadratic convergence is local and conditional; the theorem deliberately assumes the iterates remain in the regularity region.\end{proofanalysis}

\begin{computationalrelevance}This explains Newton's rapid root polishing while also documenting why global safeguarding is needed.\end{computationalrelevance}

\subsection{Optimization Algorithms}

\subsubsection{Golden-section interval contraction}

\begin{theorem}[Golden-section interval contraction]\label{prf:CONV-GOLDEN-SECTION}

For the golden-section search geometry, let
\[
\rho=\frac{\sqrt5-1}{2}.
\]
The self-similar placement of interior points makes each retained interval have length $\rho$ times its parent. Consequently, after $n$ valid reductions,
\[
L_n=\rho^n L_0\to0.
\]

\end{theorem}

\meta{Proof ID}{\pidraw{CONV-GOLDEN-SECTION}}
\meta{Object type}{CONVERGENCE-THEOREM}
\meta{Immediate dependencies}{\pidref{AX-REAL-FIELD}, \pidref{AX-ELEMENTARY-SEQUENCE-LIMITS}}
\meta{Additional transitive dependencies}{\emph{None}}
\meta{Source traceability}{\textbf{Numerical source:} Numerical Recipes, Sec. 10.2}
\meta{Catalog method IDs}{M10}

\begin{proofstrategy}Algebraic derivation followed by a geometric-sequence limit.\end{proofstrategy}

\begin{proof}

Self-similarity requires that after discarding one outer portion, one old interior
point becomes an interior point of the retained interval with the same relative
placement. If the retained fraction is $\rho$, the geometry gives
\[
1-\rho=\rho^2.
\]
The positive solution of $\rho^2+\rho-1=0$ is
\[
\rho=\frac{\sqrt5-1}{2},
\]
which satisfies $0<\rho<1$. Hence every valid step multiplies the bracket length
by $\rho$, and induction gives $L_n=\rho^nL_0$. Since $|\rho|<1$,
$L_n\to0$.

\end{proof}

\begin{proofanalysis}The theorem proves geometric contraction, not that an arbitrary objective is unimodal or that an invalid bracket contains a global minimum.\end{proofanalysis}

\begin{computationalrelevance}This is the convergence geometry behind one-dimensional golden-section minimization.\end{computationalrelevance}

\subsection{Numerical Integration}

\subsubsection{Richardson error cancellation}

\begin{proposition}[Richardson error cancellation]\label{prf:DER-RICHARDSON}

Suppose
\[
A(h)=A+c h^p+O(h^{p+q})
\]
with $p,q>0$. Then
\[
R(h)=\frac{2^pA(h/2)-A(h)}{2^p-1}
=
A+O(h^{p+q}).
\]

\end{proposition}

\meta{Proof ID}{\pidraw{DER-RICHARDSON}}
\meta{Object type}{DERIVATION}
\meta{Immediate dependencies}{\pidref{DEF-BIG-O}}
\meta{Additional transitive dependencies}{\emph{None}}
\meta{Source traceability}{\textbf{Numerical source:} Numerical Recipes, Secs. 4.3 and 17.3}
\meta{Catalog method IDs}{M15,M23}

\begin{proofstrategy}Algebraic derivation by cancellation of the leading error term.\end{proofstrategy}

\begin{proof}

From the hypothesis,
\[
A(h/2)=A+c\frac{h^p}{2^p}+O(h^{p+q}).
\]
Multiplying by $2^p$ gives
\[
2^pA(h/2)=2^pA+c h^p+O(h^{p+q}).
\]
Subtract $A(h)=A+c h^p+O(h^{p+q})$:
\[
2^pA(h/2)-A(h)=(2^p-1)A+O(h^{p+q}).
\]
Division by the nonzero constant $2^p-1$ yields the result.

\end{proof}

\begin{proofanalysis}Richardson extrapolation works because the same leading error coefficient appears at two related step sizes.\end{proofanalysis}

\begin{computationalrelevance}This identity is the core algebra behind Romberg integration and extrapolation-based ODE methods.\end{computationalrelevance}

\subsubsection{Simpson's rule is exact for polynomials of degree at most three}

\begin{theorem}[Simpson's rule is exact for polynomials of degree at most three]\label{prf:THM-SIMPSON-EXACT-CUBIC}

For any polynomial $p$ of degree at most $3$ and midpoint $m=(a+b)/2$,
\[
\int_a^b p(x)\,dx
=
\frac{b-a}{6}\left[p(a)+4p(m)+p(b)\right].
\]

\end{theorem}

\meta{Proof ID}{\pidraw{THM-SIMPSON-EXACT-CUBIC}}
\meta{Object type}{THEOREM}
\meta{Immediate dependencies}{\pidref{THM-FTC-II}}
\meta{Additional transitive dependencies}{\pidref{AX-EVT}, \pidref{AX-REAL-FIELD}, \pidref{AX-RIEMANN-BASIC}, \pidref{DEF-CONTINUITY}, \pidref{DEF-DERIVATIVE}, \pidref{DEF-RIEMANN-INTEGRAL}, \pidref{THM-FERMAT}, \pidref{THM-FTC-I}, \pidref{THM-ROLLE}, \pidref{THM-MVT}}
\meta{Source traceability}{\textbf{Mathematical source:} source calculus framework, Sec. 8.6; \textbf{Numerical source:} Numerical Recipes, Sec. 4.1}
\meta{Catalog method IDs}{M15}

\begin{proofstrategy}Direct proof by linearity and verification on a polynomial basis.\end{proofstrategy}

\begin{proof}

By linearity, it is enough to prove exactness after the affine change of variables
$x=m+ht$, where $h=(b-a)/2$, for the basis $1,t,t^2,t^3$ on $[-1,1]$.
The Simpson functional there is
\[
S(q)=\frac13\bigl(q(-1)+4q(0)+q(1)\bigr).
\]
Direct calculation gives
\[
S(1)=2=\int_{-1}^1 1\,dt,\quad
S(t)=0=\int_{-1}^1 t\,dt,
\]
\[
S(t^2)=\frac23=\int_{-1}^1 t^2\,dt,\quad
S(t^3)=0=\int_{-1}^1 t^3\,dt.
\]
Therefore the formula is exact for every cubic polynomial, and transforming back
to $[a,b]$ gives the stated rule.

\end{proof}

\begin{proofanalysis}Linearity reduces a four-dimensional polynomial statement to four elementary checks.\end{proofanalysis}

\begin{computationalrelevance}Simpson-type formulas and Romberg extrapolation build on polynomial exactness and cancellation of low-order error terms.\end{computationalrelevance}

\subsubsection{Trapezoidal rule is exact for affine functions}

\begin{theorem}[Trapezoidal rule is exact for affine functions]\label{prf:THM-TRAPEZOID-EXACT-LINEAR}

If $f(x)=mx+c$, then on $[a,b]$,
\[
\int_a^b f(x)\,dx
=
\frac{b-a}{2}\bigl(f(a)+f(b)\bigr).
\]

\end{theorem}

\meta{Proof ID}{\pidraw{THM-TRAPEZOID-EXACT-LINEAR}}
\meta{Object type}{THEOREM}
\meta{Immediate dependencies}{\pidref{THM-FTC-II}}
\meta{Additional transitive dependencies}{\pidref{AX-EVT}, \pidref{AX-REAL-FIELD}, \pidref{AX-RIEMANN-BASIC}, \pidref{DEF-CONTINUITY}, \pidref{DEF-DERIVATIVE}, \pidref{DEF-RIEMANN-INTEGRAL}, \pidref{THM-FERMAT}, \pidref{THM-FTC-I}, \pidref{THM-ROLLE}, \pidref{THM-MVT}}
\meta{Source traceability}{\textbf{Mathematical source:} source calculus framework, Sec. 8.6; \textbf{Numerical source:} Numerical Recipes, Sec. 4.1}
\meta{Catalog method IDs}{M15}

\begin{proofstrategy}Direct algebraic verification.\end{proofstrategy}

\begin{proof}

An antiderivative is $F(x)=\tfrac12mx^2+cx$. Hence
\[
\int_a^b f(x)\,dx
=
\frac{m}{2}(b^2-a^2)+c(b-a)
=
(b-a)\left(\frac{m(a+b)}2+c\right).
\]
Also
\[
\frac{f(a)+f(b)}2
=
\frac{ma+c+mb+c}{2}
=
\frac{m(a+b)}2+c.
\]
Multiplying by $b-a$ gives the same value.

\end{proof}

\begin{proofanalysis}The trapezoid formula exactly integrates the unique straight line through the endpoint values.\end{proofanalysis}

\begin{computationalrelevance}Polynomial exactness is the simplest correctness property of a quadrature rule.\end{computationalrelevance}

\subsection{Linear Algebra}

\subsubsection{Correctness of solving through an LU factorization}

\begin{theorem}[Correctness of solving through an LU factorization]\label{prf:THM-LU-SOLVE-CORRECT}

Suppose $A=LU$, where $L$ and $U$ are nonsingular triangular matrices. If $y$ solves $Ly=b$ and $x$ solves $Ux=y$, then $x$ is the unique solution of $Ax=b$.

\end{theorem}

\meta{Proof ID}{\pidraw{THM-LU-SOLVE-CORRECT}}
\meta{Object type}{THEOREM}
\meta{Immediate dependencies}{\pidref{AX-LINEAR-ALGEBRA-BASICS}}
\meta{Additional transitive dependencies}{\emph{None}}
\meta{Source traceability}{\textbf{Numerical source:} Numerical Recipes, Sec. 2.3}
\meta{Catalog method IDs}{M20}

\begin{proofstrategy}Direct proof by substitution and invertibility.\end{proofstrategy}

\begin{proof}

From $Ux=y$ and $Ly=b$,
\[
Ax=LUx=Ly=b.
\]
Thus $x$ is a solution. If $Ax_1=b=Ax_2$, then
$A(x_1-x_2)=0$. Because $A=LU$ is a product of nonsingular matrices, $A$ is
nonsingular, so $x_1-x_2=0$. Hence the solution is unique.

\end{proof}

\begin{proofanalysis}Factorization does not change the linear system; it reorganizes one solve into two triangular solves.\end{proofanalysis}

\begin{computationalrelevance}This is the algebraic correctness statement behind dense LU linear-system routines.\end{computationalrelevance}

\subsubsection{Eigenvectors of a real symmetric matrix for distinct eigenvalues are orthogonal}

\begin{theorem}[Eigenvectors of a real symmetric matrix for distinct eigenvalues are orthogonal]\label{prf:THM-SYMMETRIC-EIGENVECTORS-ORTHOGONAL}

Let $A=A^T$. If $Av=\lambda v$ and $Aw=\mu w$ with $\lambda\ne\mu$, then $v^Tw=0$.

\end{theorem}

\meta{Proof ID}{\pidraw{THM-SYMMETRIC-EIGENVECTORS-ORTHOGONAL}}
\meta{Object type}{THEOREM}
\meta{Immediate dependencies}{\pidref{AX-LINEAR-ALGEBRA-BASICS}}
\meta{Additional transitive dependencies}{\emph{None}}
\meta{Source traceability}{\textbf{Numerical source:} Numerical Recipes, Ch. 11}
\meta{Catalog method IDs}{M22}

\begin{proofstrategy}Direct algebraic proof.\end{proofstrategy}

\begin{proof}

Because $A$ is symmetric,
\[
v^TAw=(A v)^Tw.
\]
The left side is $\mu v^Tw$, while the right side is
$\lambda v^Tw$. Therefore
\[
(\mu-\lambda)v^Tw=0.
\]
Since $\mu\ne\lambda$, it follows that $v^Tw=0$.

\end{proof}

\begin{proofanalysis}Symmetry makes the two ways of moving $A$ across an inner product identical.\end{proofanalysis}

\begin{computationalrelevance}Orthogonality is a central structural property exploited by symmetric eigensystem algorithms.\end{computationalrelevance}

\subsubsection{Normal equations for full-rank linear least squares}

\begin{theorem}[Normal equations for full-rank linear least squares]\label{prf:THM-LEAST-SQUARES-NORMAL}

Let $A\in\mathbb{R}^{m\times n}$ have full column rank. A vector $x_*$ uniquely minimizes
$\|Ax-b\|_2^2$ if and only if
\[
A^TAx_*=A^Tb.
\]

\end{theorem}

\meta{Proof ID}{\pidraw{THM-LEAST-SQUARES-NORMAL}}
\meta{Object type}{THEOREM}
\meta{Immediate dependencies}{\pidref{AX-LINEAR-ALGEBRA-BASICS}, \pidref{THM-FERMAT}}
\meta{Additional transitive dependencies}{\pidref{AX-REAL-FIELD}, \pidref{DEF-DERIVATIVE}}
\meta{Source traceability}{\textbf{Numerical source:} Numerical Recipes, Sec. 15.4 and Ch. 2}
\meta{Catalog method IDs}{M21,M29}

\begin{proofstrategy}Direct proof by orthogonality and quadratic expansion.\end{proofstrategy}

\begin{proof}

Suppose first that $A^T(Ax_*-b)=0$. For any $h$,
\[
\|A(x_*+h)-b\|_2^2
=
\|Ax_*-b\|_2^2
+2h^TA^T(Ax_*-b)+\|Ah\|_2^2.
\]
The middle term is zero, so the objective cannot decrease. Because $A$ has full
column rank, $Ah=0$ only when $h=0$, hence the minimizer is unique.

Conversely, if $x_*$ minimizes the objective, fix any vector $v$ and consider
\[
\psi(t)=\|A(x_*+tv)-b\|_2^2.
\]
The scalar function has a minimum at $t=0$, so by \pid{THM-FERMAT},
\[
0=\psi'(0)=2v^TA^T(Ax_*-b).
\]
This holds for every $v$, forcing $A^T(Ax_*-b)=0$, which is the normal equation.

\end{proof}

\begin{proofanalysis}The residual at a least-squares solution is orthogonal to every column of the design matrix.\end{proofanalysis}

\begin{computationalrelevance}This theorem is the mathematical core of linear regression and parameter fitting before numerical conditioning motivates QR or SVD implementations.\end{computationalrelevance}

\subsection{Differential Equations}

\subsubsection{Local truncation error of Euler's method}

\begin{theorem}[Local truncation error of Euler's method]\label{prf:ERR-EULER-LOCAL}

If the exact solution $y$ of $y'=f(t,y)$ has a continuous second derivative, then one Euler step from the exact value satisfies
\[
y(t+h)=y(t)+h f(t,y(t))+\tau,\qquad
|\tau|\le \frac{h^2}{2}\max_{s\text{ between }t\text{ and }t+h}|y''(s)|.
\]
Thus Euler's one-step local truncation error is $O(h^2)$.

\end{theorem}

\meta{Proof ID}{\pidraw{ERR-EULER-LOCAL}}
\meta{Object type}{ERROR-THEOREM}
\meta{Immediate dependencies}{\pidref{THM-TAYLOR-LAGRANGE}, \pidref{DEF-BIG-O}}
\meta{Additional transitive dependencies}{\pidref{AX-EVT}, \pidref{AX-REAL-FIELD}, \pidref{DEF-DERIVATIVE}, \pidref{THM-FERMAT}, \pidref{THM-ROLLE}}
\meta{Source traceability}{\textbf{Mathematical source:} source calculus framework, Sec. 6.1; \textbf{Numerical source:} Numerical Recipes, Ch. 17}
\meta{Catalog method IDs}{M23}

\begin{proofstrategy}Direct Taylor expansion of the exact solution.\end{proofstrategy}

\begin{proof}

Taylor's theorem gives
\[
y(t+h)=y(t)+hy'(t)+\frac{h^2}{2}y''(\xi)
\]
for some $\xi$ between $t$ and $t+h$. Since the differential equation gives
$y'(t)=f(t,y(t))$,
\[
y(t+h)=y(t)+h f(t,y(t))+\frac{h^2}{2}y''(\xi).
\]
The bound and $O(h^2)$ statement follow immediately.

\end{proof}

\begin{proofanalysis}Consistency of Euler's method is a Taylor statement about one exact step; global convergence requires an additional stability argument not claimed here.\end{proofanalysis}

\begin{computationalrelevance}This separates local method order from the stronger global-error guarantees of production ODE solvers.\end{computationalrelevance}

\subsection{Spectral Methods}

\subsubsection{Radix-2 even/odd decomposition of the DFT}

\begin{theorem}[Radix-2 even/odd decomposition of the DFT]\label{prf:THM-DFT-EVEN-ODD}

Let $N=2m$, $\omega_N=e^{-2\pi i/N}$, and
\[
X_k=\sum_{n=0}^{N-1}x_n\omega_N^{kn}.
\]
Then
\[
X_k=E_k+\omega_N^k O_k,
\]
where
\[
E_k=\sum_{j=0}^{m-1}x_{2j}\omega_m^{kj},
\qquad
O_k=\sum_{j=0}^{m-1}x_{2j+1}\omega_m^{kj}.
\]

\end{theorem}

\meta{Proof ID}{\pidraw{THM-DFT-EVEN-ODD}}
\meta{Object type}{THEOREM}
\meta{Immediate dependencies}{\pidref{AX-COMPLEX-EXP}}
\meta{Additional transitive dependencies}{\emph{None}}
\meta{Source traceability}{\textbf{Numerical source:} Numerical Recipes, Ch. 12}
\meta{Catalog method IDs}{M26,M27}

\begin{proofstrategy}Algebraic decomposition by parity.\end{proofstrategy}

\begin{proof}

Split the DFT sum into even and odd indices:
\[
X_k=
\sum_{j=0}^{m-1}x_{2j}\omega_N^{k(2j)}
+
\sum_{j=0}^{m-1}x_{2j+1}\omega_N^{k(2j+1)}.
\]
Because $\omega_N^2=\omega_m$,
\[
X_k=
\sum_{j=0}^{m-1}x_{2j}\omega_m^{kj}
+
\omega_N^k\sum_{j=0}^{m-1}x_{2j+1}\omega_m^{kj}
=
E_k+\omega_N^kO_k.
\]

\end{proof}

\begin{proofanalysis}The identity turns one size-$N$ transform into two size-$N/2$ transforms plus linear-time recombination.\end{proofanalysis}

\begin{computationalrelevance}This is the algebraic factorization at the heart of radix-2 FFT implementations.\end{computationalrelevance}

\subsubsection{Radix-2 FFT has \$O(N\textbackslash{}log N)\$ arithmetic complexity}

\begin{theorem}[Radix-2 FFT has \$O(N\textbackslash{}log N)\$ arithmetic complexity]\label{prf:THM-FFT-COMPLEXITY}

If a radix-2 FFT on $N=2^m$ samples satisfies
\[
T(N)\le 2T(N/2)+cN
\]
for $N\ge2$ and fixed $c$, then $T(N)=O(N\log_2N)$.

\end{theorem}

\meta{Proof ID}{\pidraw{THM-FFT-COMPLEXITY}}
\meta{Object type}{THEOREM}
\meta{Immediate dependencies}{\pidref{THM-DFT-EVEN-ODD}, \pidref{DEF-BIG-O}}
\meta{Additional transitive dependencies}{\pidref{AX-COMPLEX-EXP}}
\meta{Source traceability}{\textbf{Numerical source:} Numerical Recipes, Ch. 12}
\meta{Catalog method IDs}{M26}

\begin{proofstrategy}Mathematical induction on \$m=\textbackslash{}log\_2N\$.\end{proofstrategy}

\begin{proof}

Let $N=2^m$. We prove by induction that
\[
T(2^m)\le C\,2^m m+D\,2^m
\]
for suitable constants $C\ge c$ and $D$ large enough for the base case.
Assume the inequality for $m-1$. Then
\[
T(2^m)
\le
2T(2^{m-1})+c2^m
\]
\[
\le
2\left(C2^{m-1}(m-1)+D2^{m-1}\right)+c2^m
=
C2^m(m-1)+D2^m+c2^m.
\]
Since $C\ge c$, this is at most
$C2^m m+D2^m$. Therefore
$T(N)\le CN\log_2N+DN$, which is $O(N\log N)$.

\end{proof}

\begin{proofanalysis}The complexity reduction follows from the recursive factorization, not from an approximation of the DFT.\end{proofanalysis}

\begin{computationalrelevance}It explains why spectral methods become practical for large sampled signals and grids.\end{computationalrelevance}

\subsection{Probability and Monte Carlo}

\subsubsection{Variance of the independent Monte Carlo mean}

\begin{theorem}[Variance of the independent Monte Carlo mean]\label{prf:ERR-MC-VARIANCE}

If $X_1,\ldots,X_N$ are independent with common finite variance $\sigma^2$, then
\[
\operatorname{Var}(\widehat\mu_N)=\frac{\sigma^2}{N},
\qquad
\operatorname{sd}(\widehat\mu_N)=\frac{\sigma}{\sqrt N}.
\]

\end{theorem}

\meta{Proof ID}{\pidraw{ERR-MC-VARIANCE}}
\meta{Object type}{ERROR-THEOREM}
\meta{Immediate dependencies}{\pidref{AX-PROBABILITY}}
\meta{Additional transitive dependencies}{\emph{None}}
\meta{Source traceability}{\textbf{Numerical source:} Numerical Recipes, Secs. 7.7--7.9}
\meta{Catalog method IDs}{M18}

\begin{proofstrategy}Direct variance calculation.\end{proofstrategy}

\begin{proof}

Using independence and scaling of variance,
\[
\operatorname{Var}(\widehat\mu_N)
=
\operatorname{Var}\left(\frac1N\sum_{i=1}^N X_i\right)
=
\frac1{N^2}\sum_{i=1}^N\operatorname{Var}(X_i)
=
\frac{N\sigma^2}{N^2}
=
\frac{\sigma^2}{N}.
\]
Taking square roots gives the standard-deviation formula.

\end{proof}

\begin{proofanalysis}The characteristic $N^{-1/2}$ Monte Carlo error scale is a variance statement, not a deterministic quadrature order.\end{proofanalysis}

\begin{computationalrelevance}It explains why naive Monte Carlo converges slowly but dimension-independently compared with grid quadrature.\end{computationalrelevance}

\subsubsection{Monte Carlo sample mean is unbiased}

\begin{theorem}[Monte Carlo sample mean is unbiased]\label{prf:THM-MC-UNBIASED}

Let $X_1,\ldots,X_N$ have common finite expectation $\mu$. Then
\[
\widehat\mu_N=\frac1N\sum_{i=1}^N X_i
\]
satisfies $E[\widehat\mu_N]=\mu$.

\end{theorem}

\meta{Proof ID}{\pidraw{THM-MC-UNBIASED}}
\meta{Object type}{THEOREM}
\meta{Immediate dependencies}{\pidref{AX-PROBABILITY}}
\meta{Additional transitive dependencies}{\emph{None}}
\meta{Source traceability}{\textbf{Numerical source:} Numerical Recipes, Sec. 7.7}
\meta{Catalog method IDs}{M18,M28}

\begin{proofstrategy}Direct proof by linearity of expectation.\end{proofstrategy}

\begin{proof}

By linearity of expectation,
\[
E[\widehat\mu_N]
=
\frac1N\sum_{i=1}^N E[X_i]
=
\frac1N\sum_{i=1}^N\mu
=
\mu.
\]
Independence is not required for unbiasedness.

\end{proof}

\begin{proofanalysis}Averaging does not introduce bias when every sample has the correct expectation.\end{proofanalysis}

\begin{computationalrelevance}Simple Monte Carlo integration estimates an integral as an expectation and then replaces that expectation by a sample mean.\end{computationalrelevance}

\subsection{PDE Discretization}

\subsubsection{Second-order centered approximation to a second derivative}

\begin{theorem}[Second-order centered approximation to a second derivative]\label{prf:ERR-SECOND-DERIVATIVE-CENTERED}

If $u$ has a continuous fourth derivative on $[x-h,x+h]$, then
\[
\frac{u(x+h)-2u(x)+u(x-h)}{h^2}
=
u''(x)+E_h,
\]
with
\[
|E_h|
\le
\frac{h^2}{12}
\max_{|t-x|\le |h|}|u^{(4)}(t)|.
\]

\end{theorem}

\meta{Proof ID}{\pidraw{ERR-SECOND-DERIVATIVE-CENTERED}}
\meta{Object type}{ERROR-THEOREM}
\meta{Immediate dependencies}{\pidref{THM-TAYLOR-LAGRANGE}}
\meta{Additional transitive dependencies}{\pidref{AX-EVT}, \pidref{AX-REAL-FIELD}, \pidref{DEF-DERIVATIVE}, \pidref{THM-FERMAT}, \pidref{THM-ROLLE}}
\meta{Source traceability}{\textbf{Numerical source:} Numerical Recipes, Ch. 20}
\meta{Catalog method IDs}{M33}

\begin{proofstrategy}Direct proof by symmetric fourth-order Taylor expansions.\end{proofstrategy}

\begin{proof}

Taylor's theorem gives
\[
u(x\pm h)
=
u(x)\pm hu'(x)+\frac{h^2}{2}u''(x)
\pm\frac{h^3}{6}u'''(x)
+\frac{h^4}{24}u^{(4)}(\xi_\pm).
\]
Adding the two formulas and subtracting $2u(x)$ cancels the odd terms:
\[
u(x+h)-2u(x)+u(x-h)
=
h^2u''(x)
+\frac{h^4}{24}\left(u^{(4)}(\xi_+)+u^{(4)}(\xi_-)\right).
\]
Divide by $h^2$ and bound the two fourth derivatives by their maximum to obtain
the stated estimate.

\end{proof}

\begin{proofanalysis}The centered stencil gains second-order accuracy because odd Taylor terms cancel.\end{proofanalysis}

\begin{computationalrelevance}This is the canonical local consistency result behind finite-difference discretizations of diffusion, Poisson, and wave operators.\end{computationalrelevance}
\begin{landscape}
\section{Dependency Expansion Audit}
\small
\setlength{\tabcolsep}{8pt}
\begin{longtable}{@{}lrrrr@{}}
\toprule
\textbf{Proof ID} & \textbf{Immediate} & \textbf{Transitive} & \textbf{Expanded} & \textbf{Foundation reached}\\
\midrule
\endfirsthead
\toprule
\textbf{Proof ID} & \textbf{Immediate} & \textbf{Transitive} & \textbf{Expanded} & \textbf{Foundation reached}\\
\midrule
\endhead
\pidref{THM-DFT-EVEN-ODD} & 1 & 1 & Yes & Yes \\
\pidref{THM-LU-SOLVE-CORRECT} & 1 & 1 & Yes & Yes \\
\pidref{THM-SYMMETRIC-EIGENVECTORS-ORTHOGONAL} & 1 & 1 & Yes & Yes \\
\pidref{LEM-POLYNOMIAL-ROOT-COUNT} & 1 & 1 & Yes & Yes \\
\pidref{ERR-MC-VARIANCE} & 1 & 1 & Yes & Yes \\
\pidref{THM-MC-UNBIASED} & 1 & 1 & Yes & Yes \\
\pidref{CONV-GOLDEN-SECTION} & 2 & 2 & Yes & Yes \\
\pidref{THM-NESTED-INTERVAL} & 2 & 2 & Yes & Yes \\
\pidref{DER-RICHARDSON} & 1 & 1 & Yes & No \\
\pidref{THM-IVT} & 3 & 3 & Yes & Yes \\
\pidref{DER-NEWTON-ITERATION} & 1 & 1 & Yes & No \\
\pidref{THM-FERMAT} & 2 & 2 & Yes & Yes \\
\pidref{THM-CURVE-CHAIN} & 1 & 1 & Yes & No \\
\pidref{THM-DIRECTIONAL-GRADIENT} & 2 & 2 & Yes & No \\
\pidref{LEM-LIMIT-BOUNDED} & 2 & 2 & Yes & Yes \\
\pidref{LEM-RECIPROCAL-BOUND} & 2 & 2 & Yes & Yes \\
\pidref{THM-LIMIT-SUM} & 1 & 1 & Yes & No \\
\pidref{THM-LIMIT-UNIQUENESS} & 2 & 2 & Yes & Yes \\
\pidref{THM-FTC-I} & 3 & 3 & Yes & Yes \\
\pidref{THM-GEOMETRIC-SERIES} & 3 & 3 & Yes & Yes \\
\pidref{THM-FFT-COMPLEXITY} & 2 & 3 & Yes & Yes \\
\pidref{THM-INTERPOLATION-UNIQUENESS} & 1 & 2 & Yes & Yes \\
\pidref{ALG-SIGN-BRACKET-INVARIANT} & 1 & 4 & Yes & Yes \\
\pidref{THM-LEAST-SQUARES-NORMAL} & 2 & 4 & Yes & Yes \\
\pidref{THM-ROLLE} & 2 & 4 & Yes & Yes \\
\pidref{THM-LIMIT-PRODUCT} & 3 & 4 & Yes & Yes \\
\pidref{THM-SERIES-TERM-TEST} & 2 & 3 & Yes & No \\
\pidref{THM-LAGRANGE-INTERPOLATION} & 1 & 3 & Yes & Yes \\
\pidref{ALG-BISECTION-INVARIANT} & 1 & 5 & Yes & Yes \\
\pidref{THM-MVT} & 1 & 5 & Yes & Yes \\
\pidref{THM-TAYLOR-LAGRANGE} & 1 & 5 & Yes & Yes \\
\pidref{THM-DIFFERENTIABLE-CONTINUOUS} & 3 & 7 & Yes & Yes \\
\pidref{THM-LIMIT-QUOTIENT} & 3 & 6 & Yes & Yes \\
\pidref{CONV-BISECTION} & 3 & 7 & Yes & Yes \\
\pidref{COR-MONOTONICITY} & 1 & 6 & Yes & Yes \\
\pidref{THM-FTC-II} & 2 & 10 & Yes & Yes \\
\pidref{CONV-NEWTON-LOCAL} & 2 & 7 & Yes & Yes \\
\pidref{ERR-CENTERED-DIFFERENCE} & 2 & 7 & Yes & Yes \\
\pidref{ERR-EULER-LOCAL} & 2 & 7 & Yes & Yes \\
\pidref{ERR-FORWARD-DIFFERENCE} & 2 & 7 & Yes & Yes \\
\pidref{ERR-SECOND-DERIVATIVE-CENTERED} & 1 & 6 & Yes & Yes \\
\pidref{THM-CHAIN-RULE} & 3 & 8 & Yes & Yes \\
\pidref{THM-PRODUCT-RULE} & 4 & 8 & Yes & Yes \\
\pidref{THM-SECOND-DERIVATIVE-TEST} & 2 & 7 & Yes & Yes \\
\pidref{THM-FUNDAMENTAL-LINE-INTEGRAL} & 2 & 13 & Yes & Yes \\
\pidref{THM-SIMPSON-EXACT-CUBIC} & 1 & 11 & Yes & Yes \\
\pidref{THM-TRAPEZOID-EXACT-LINEAR} & 1 & 11 & Yes & Yes \\
\pidref{COR-IMPLICIT-DERIVATIVE} & 2 & 10 & Yes & Yes \\
\pidref{DER-PARAMETRIC-DERIVATIVE} & 1 & 9 & Yes & Yes \\
\pidref{DER-SUBSTITUTION} & 3 & 17 & Yes & Yes \\
\pidref{DER-INTEGRATION-BY-PARTS} & 2 & 17 & Yes & Yes \\
\pidref{THM-QUOTIENT-RULE} & 4 & 11 & Yes & Yes \\
\pidref{COR-PATH-INDEPENDENCE} & 1 & 14 & Yes & Yes \\
\bottomrule
\end{longtable}
\end{landscape}

\section{References and Source Roles}

\begin{thebibliography}{9}

\bibitem{proofs}
Gary Chartrand, Albert D. Polimeni, and Ping Zhang,
\emph{Mathematical Proofs: A Transition to Advanced Mathematics},
Fourth Edition, Pearson, 2018.
This is the proof-writing and proof-methodology authority for the library.

\bibitem{nr}
William H. Press, Saul A. Teukolsky, William T. Vetterling, and Brian P. Flannery,
\emph{Numerical Recipes: The Art of Scientific Computing},
Third Edition, Cambridge University Press, 2007.
This is the numerical-method taxonomy and numerical-analysis context source.

\bibitem{calcsource}
\emph{Source calculus framework}.
The supplied calculus source provides definitions, theorem statements, instructional ordering,
and application context. It is intentionally not identified by title or author in this document.

\end{thebibliography}

\section*{Validation Summary}
The generated registry is acyclic. Every explicit dependency ID exists. The modular and complete
documents are generated from the same canonical statement registry, so their theorem statements,
hypotheses, conclusions, proof IDs, and source mappings are mechanically identical. The complete
version differs only in dependency expansion depth.

\end{document}
